\documentclass[11pt]{amsart}

% Submission build for the Journal of the London Mathematical Society.
% JLMS requests a single PDF at initial submission and does not prescribe a class.

% Packages
\usepackage[T1]{fontenc}
\usepackage{newtxtext}
\usepackage{newtxmath}
\usepackage{textcomp}
\usepackage{microtype}
\usepackage{graphicx}
\usepackage{booktabs}
\usepackage{array}
\usepackage{enumitem}
\usepackage[numbers,sort&compress]{natbib}
\usepackage{xcolor}
\definecolor{linkblue}{RGB}{25,70,130}
\usepackage[colorlinks=true,linkcolor=linkblue,citecolor=linkblue,urlcolor=linkblue]{hyperref}
\usepackage{aliascnt}
\usepackage[nameinlink,capitalise,noabbrev]{cleveref}
\raggedbottom

% Theorem environments: one shared counter per section.
\newtheorem{theorem}{Theorem}[section]
\newaliascnt{lemma}{theorem}
\newtheorem{lemma}[lemma]{Lemma}
\aliascntresetthe{lemma}
\crefname{lemma}{Lemma}{Lemmas}
\newaliascnt{proposition}{theorem}
\newtheorem{proposition}[proposition]{Proposition}
\aliascntresetthe{proposition}
\crefname{proposition}{Proposition}{Propositions}
\newaliascnt{corollary}{theorem}
\newtheorem{corollary}[corollary]{Corollary}
\aliascntresetthe{corollary}
\crefname{corollary}{Corollary}{Corollaries}
\theoremstyle{definition}
\newaliascnt{definition}{theorem}
\newtheorem{definition}[definition]{Definition}
\aliascntresetthe{definition}
\crefname{definition}{Definition}{Definitions}
\theoremstyle{remark}
\newaliascnt{remark}{theorem}
\newtheorem{remark}[remark]{Remark}
\aliascntresetthe{remark}
\crefname{remark}{Remark}{Remarks}
% Lettered main theorems (Section 1)
\theoremstyle{plain}
\newtheorem{lettertheorem}{Theorem}
\renewcommand{\thelettertheorem}{\Alph{lettertheorem}}
\crefname{lettertheorem}{Theorem}{Theorems}
\numberwithin{equation}{section}

% Math operators / commands
\DeclareMathOperator{\tr}{tr}
\DeclareMathOperator{\Id}{Id}
\DeclareMathOperator{\negl}{negl}
\DeclareMathOperator{\Ind}{Ind}
\DeclareMathOperator{\GL}{GL}
\DeclareMathOperator{\SL}{SL}
\DeclareMathOperator{\PGL}{PGL}
\DeclareMathOperator{\PSL}{PSL}
\DeclareMathOperator{\Sym}{Sym}
\DeclareMathOperator{\Stab}{Stab}
\DeclareMathOperator{\Fix}{Fix}
\newcommand{\norm}[1]{\lVert #1 \rVert}
\newcommand{\abs}[1]{\lvert #1 \rvert}
\newcommand{\frob}[1]{\norm{#1}_{\mathrm{F}}}
\newcommand{\opn}[1]{\norm{#1}}
\newcommand{\dtv}{d_{\mathrm{TV}}}
\newcommand{\BP}{\mathrm{BP}}
\newcommand{\F}{\mathbb{F}}
\newcommand{\Z}{\mathbb{Z}}
\newcommand{\C}{\mathbb{C}}
\newcommand{\R}{\mathbb{R}}

\title[The solvability dichotomy for commutator-tree mixing]
{The solvability dichotomy for commutator-tree mixing in finite groups}
\author{Ernest Darell Zermeño}
\address{Universidad Panamericana, Campus Guadalajara, Mexico}
\email{0244552@up.edu.mx}
\subjclass[2020]{Primary 20F14; Secondary 20C15, 20D10, 20P05, 60B15}
\keywords{finite groups, commutator trees, solvability, derived length, conjugacy classes,
random walks on groups}
\hypersetup{
  pdftitle={The solvability dichotomy for commutator-tree mixing in finite groups},
  pdfauthor={Ernest Darell Zermeño},
  pdfsubject={Finite groups, commutator trees, and solvability},
  pdfkeywords={finite groups, commutator trees, solvability, derived length, conjugacy classes,
    random walks on groups}
}

\begin{document}

\begin{abstract}
We study random products prescribed by balanced commutator trees, with
independently configurable leaf labels and independent choices at every
expanded occurrence. A tree is valid if its deterministic commutator
value is non-identity. For a finite non-solvable group and a
conjugation-closed generating leaf set, a uniform strict-contraction
level exists exactly when every proper leaf-generated subgroup is
solvable. Its exact value is the largest derived length of these
subgroups. Among finite non-abelian simple groups, this property for
every non-trivial conjugacy class characterises the minimal simple
groups. We classify projective rank-one involution classes and their
obstructions. An arbitrary-order stability estimate and bounded-rank
diameter bounds give fixed-error mixing in optimal depth
$\Theta(\log\log|G|)$ for all valid trees in the unobstructed rank-one
families and all minimal simple groups. For primes
$q\equiv\pm3\pmod8$, $q\ne3$, every pair of exterior $S_4$ subgroups
of $\PGL_2(q)$ with intersection of order six is simultaneously
conjugate to a natural Hamilton pair. Compression of a quaternionic
Ramanujan walk bounds the largest fixed-space cosine over non-trivial
irreducibles by $\sqrt3/2$, with asymptotic equality.
\end{abstract}

\maketitle



%% ===================================================================
\section{Introduction}
\label{sec:intro}
%% ===================================================================

Random products in a finite group provide a basic way to approach the uniform
distribution. Each step may simply multiply by a prescribed group element or
do nothing, with equal probability. We study what happens when the order of
these steps is prescribed by repeatedly forming commutators
$[x,y]=xyx^{-1}y^{-1}$. We require the underlying deterministic word to have
non-identity value. Under this constraint, when does every allowed configuration
begin to contract towards the uniform distribution, and what determines the
required depth?

The answer is controlled by solvability. Iterated commutators in a solvable
subgroup eventually become trivial, at a depth measured by its derived length.
The non-identity requirement therefore prevents sufficiently deep configurations
from remaining inside such a subgroup. We prove that, for a conjugation-closed
set of allowed elements generating a finite non-solvable group, the exact
contraction threshold is the largest derived length of a proper subgroup
generated by allowed elements. If one of those proper subgroups is non-solvable,
it instead supplies an obstruction at every depth.

Equivalently, we ask when a non-identity iterated commutator certifies that
its entries generate the whole group. Once the depth reaches the stated
threshold, no proper subgroup generated by allowed entries can support such
a value. This gives a group-theoretic interpretation of the contraction law.

Transpositions illustrate the distinction. In $S_5$, every proper subgroup
generated by transpositions is solvable, and the largest derived length is
$3$. In $S_n$ for $n\ge6$, the point stabilizer $S_{n-1}$ is already non-solvable:
at every depth it supplies valid configurations supported in a proper
subgroup. Thus uniform contraction holds for $S_5$ and is obstructed for
every $n\ge6$ (\Cref{cor:Sn}).

\paragraph{The model and the structural results.}
The leaves of a balanced binary tree are labelled by elements of a set
$\mathcal C\subseteq G$, and each internal vertex takes the commutator of its
two child values. A configuration is \emph{valid} when the root value is
non-identity. In the random product obtained by expanding the tree, each
occurrence is read independently: it contributes its prescribed element or
the identity with equal probability. Thus repeated occurrences have fresh
random choices; this is the \emph{independent-read} model. The labels at
different leaf positions may also be chosen independently, which we call
the \emph{configurable} model. Precise definitions are in \Cref{sec:prelim}.

A subgroup is $\mathcal C$-generated if it is generated by its intersection
with $\mathcal C$. Strict contraction means that the Fourier operator has
norm less than one in every non-trivial irreducible representation; a uniform
level requires this for every valid configuration at that level. The first
main theorem identifies this analytic threshold with a subgroup invariant.

\begin{lettertheorem}[Solvability dichotomy and exact level]
\label{ithm:A}
Let $G$ be a finite non-solvable group and $\mathcal{C}$ a configurable conjugation-closed
leaf set with $\langle\mathcal{C}\rangle=G$. The construction has a uniform
strict-contraction level if and only if every proper $\mathcal{C}$-generated subgroup of
$G$ is solvable; if some proper $\mathcal{C}$-generated subgroup is non-solvable, the
obstruction is permanent: every depth admits a valid configuration that fails to contract in some non-trivial irreducible. When the level exists, it
is exactly
$D=\max\{\operatorname{dl}(H):H<G,\ H\text{ is }\mathcal{C}\text{-generated}\}$, the
largest derived length among proper $\mathcal{C}$-generated subgroups
\textup{(\Cref{thm:dichotomy})}.
\end{lettertheorem}

The analytic ingredient concerns one configuration at a time. For a tree
$\tau$, let $\mu_\tau$ be its unconditional output law and write
$M_\tau(\rho)=\widehat{\mu_\tau}(\rho)$ for its Fourier operator.
The subspace fixed by $H$ is denoted by $V_\rho^H$.
Throughout, norms of operators are operator norms, written $\norm{\cdot}$.

\begin{lettertheorem}[Fourier fixed-vector criterion]
\label{ithm:B}
For any unitary representation $\rho$ and any configuration $\tau$ with leaf set $S$ and
$H=\langle S\rangle$, one has $\opn{M_\tau(\rho)}=1\iff V_\rho^H\neq0$. Hence the
configuration contracts strictly in every non-trivial irreducible if and only if
$\langle S\rangle=G$: strict contraction $\iff$ generation $\iff$ escape from every proper
subgroup \textup{(\Cref{lem:geom,thm:criterion})}.
\end{lettertheorem}

The common-fixed-vector argument proves \Cref{ithm:B}. The normal-subset
word-value lemma of Fern\'andez-Alcober and
Pintonello~\cite[Lemma 2.1(ii)]{FernandezPintonello2023} identifies the subgroup
generated by depth-$k$ targets with $H^{(k)}$. Together they identify the
exact depth at which a non-identity target forces generation, proving
\Cref{ithm:A}. \Cref{sec:notfolklore} summarises this correspondence.

For involution leaves the averaging factors are orthogonal projections onto
fixed spaces, connecting the criterion with the theory of products of
projections. For leaves of arbitrary order, equality in the contraction
inequality still characterises the fixed space.
Angles between two subgroup-fixed spaces will govern the arithmetic
application below.

Testing all non-trivial irreducibles matters: a proper subgroup can act
transitively on a chosen permutation set, so contraction on its augmentation
representation alone need not detect the obstruction. \Cref{prop:std} explains
this distinction; $S_5$ acting on six points supplies explicit examples.

For \emph{rigid} templates---one fixed leaf per position, the regime of an actual
branching program---the configurable level $D$ becomes an upper bound, and persistence of
validity is a finite dynamical question on conjugacy classes, which we solve.

\begin{lettertheorem}[Rigid templates]
\label{ithm:C}
Let $\mathcal{C}$ be conjugation-closed.
\textup{(i)} For a non-empty set $S$ of non-trivial conjugacy classes, valid binary
compositions persist at every level from $S$ if and only if some class in $S$ has
non-solvable normal closure---equivalently, escapes the solvable radical
$O_\infty(G)$---if and only if some class in $S$ reaches a cycle of the target transition
graph $\mathcal{T}(G)$; and mixed-class children never extend the reach of persistence
beyond conjugated doubling \textup{(\Cref{thm:diagred,cor:binaryred,cor:binarylevel})}.
\textup{(ii)} Over the transposition class of $S_5$, the non-vacuously realisable exact
levels of rigid families are precisely $\{2,3\}$ \textup{(\Cref{thm:s5levels})}.
\end{lettertheorem}

\paragraph{Classification and examples.}
The subgroup criterion yields exact levels in concrete families. For
involution leaves in $A_5$ the level is $2$: although $A_4$ is a proper
subgroup of derived length $2$, its involutions generate only $V_4$, so
only the subgroups generated by the allowed leaves enter the criterion
(\Cref{cor:A5}). The same criterion also characterises a classical family:
a finite non-abelian simple group is minimal simple if and only if every
non-trivial conjugacy class has a uniform strict-contraction level
(\Cref{cor:minsimple}). In that case every class has level at most $5$, and
the involution class of $\PSL_3(3)$ attains $5$
(\Cref{prop:psl33-level}). The converse uses the solvable radical of a
proper non-solvable subgroup; the bound uses the known bound on
proper-subgroup derived lengths~\cite[Remark 2.3]{ErsoyTortoraTota2014}.
The transvection class in the same group has level $4$:
\Cref{prop:transvection-levels} classifies transvection leaves in
$\SL_n(q)$ for $n\ge3$, making the dependence on the leaf set explicit.
For a generating conjugacy class in $\GL_2(\Z/p^r)$, $p\ge5$,
\Cref{prop:congruence-class} transfers the dichotomy from reduction modulo
$p$; at $r=2$ the level increases by at most one.

In projective rank one the criterion gives a full involution-class
classification. Here an \emph{exterior} involution of $\PGL_2(q)$ is one
whose determinant is a nonsquare; the other involutions lie in $\PSL_2(q)$.

\begin{lettertheorem}[Rank-one involutions]
\label{ithm:D}
For odd $q\ge5$, all involutions of $\PGL_2(q)$ are obstructed, whereas the
exterior class is obstructed exactly by proper subfields of size at least $5$
and odd extension degree. In the unobstructed exterior cases the exact level
is $3$ for $q\equiv\pm3\pmod8$ and $2$ for $q\equiv\pm1\pmod8$.
The involution classes of $\PSL_2(q)$, including its small exceptions, are
classified in \Cref{thm:rankone-classification}.
\end{lettertheorem}

\paragraph{Quantitative consequences and an arithmetic application.}
For a fixed group, strict contraction gives a positive gap by finiteness.
Over a growing family, that gap may tend to zero. For involution leaves and
uniformly bounded depth, a uniform contraction gap is equivalent to a
uniform spectral gap for the associated lazy Cayley walks: each walk chooses
an expanded leaf occurrence uniformly and applies it with probability
$1/2$. \Cref{prop:cayley-comparison,cor:bounded-depth-gap} give explicit
two-sided comparisons. This connects the block problem to expansion of its
particular generating sets.

The next result controls the mixing depth even when the base gap decreases
with group size. Write $U_G$ for the uniform law on $G$.

\begin{lettertheorem}[Uniform rates and optimal order of depth]
\label{ithm:E}
For the unobstructed classes of \Cref{ithm:D} and
\Cref{thm:rankone-classification}, and for every non-trivial conjugacy class
of every minimal simple group, there are absolute $c,B>0$ such that every
valid tree $\tau$ of level $m\ge D$ satisfies
\[
 \dtv(\mu_\tau,U_G)\le\tfrac12\sqrt{|G|-1}
                  \exp\bigl(-c4^m/(\log|G|)^B\bigr).
\]
For fixed error in $(0,1/2)$ the least depth guaranteeing this error for
all valid trees has order $\Theta(\log\log|G|)$ as $|G|\to\infty$
\textup{(\Cref{thm:family-rate,prop:depth-lower,cor:family-examples})}.
\end{lettertheorem}

The proof converts the bounded-rank diameter theorem of
Pyber--Szab\'o~\cite{PyberSzabo2016} into a base-block contraction estimate
and amplifies it through the tree. An elementary support bound gives the
matching order of depth. Over prime fields outside a set of relative density
zero among the primes, strong uniform expansion gives a constant gap for
every generating involutive block of bounded depth. For the unobstructed
classes this improves the exponential factor to $e^{-c4^m}$
(\Cref{cor:almost-all-primes}).
Over arbitrary finite fields, a second consequence of uniform expansion
bounds the number of exceptional conjugacy classes of generating
involutive configurations at each fixed depth by $O(q^\delta)$ for every
$\delta>0$ (\Cref{cor:few-exceptional-blocks}).

The connection with the structural theory is the geometry of fixed spaces:
once a common fixed vector is absent, their angle measures how strongly
averaging over two subgroups contracts. For the exterior octahedral pairs
with intersection of order six, all simultaneously conjugate to a natural
Hamilton pair, this angle admits an asymptotically optimal uniform bound.

\begin{lettertheorem}[The natural octahedral pair]
\label{ithm:F}
For primes $q\equiv\pm3\pmod8$, $q\ne3$, the Hamilton construction
defines exterior subgroups $A_q,B_q\cong S_4$ of $G_q=\PGL_2(q)$ with
$A_q\cap B_q\cong S_3$ and $\langle A_q,B_q\rangle=G_q$. Every pair
$A_1,A_2$ of exterior octahedral subgroups of $G_q$ with $|A_1\cap A_2|=6$
is simultaneously conjugate to $(A_q,B_q)$, and its fixed-space
projections satisfy
\[
 \max_{\rho\in\widehat{G_q},\,\rho\ne\mathbf1}
              \opn{P_{A_1}^\rho P_{A_2}^\rho}=c_q\le\frac{\sqrt3}{2},
 \qquad c_q\longrightarrow\frac{\sqrt3}{2},
\]
where $c_q$ is the common value attained by the natural pair
\textup{(\Cref{prop:octahedral-cap6,thm:octahedral-sharp})}.
\end{lettertheorem}

The compression identity transfers the Ramanujan bound of
Davidoff--Sarnak--Valette~\cite[Theorem 4.2.2]{DSV2003} from the norm-$3$
quaternionic walk to the angle operator. Alon--Boppana on the associated
Hecke graphs gives optimality. Geometrically,
$\theta_q=\arccos c_q$ is the smallest fixed-space angle over all non-trivial
irreducibles: $\theta_q\ge\pi/6$ and $\theta_q\to\pi/6$. Combined with an
internal $S_4$ estimate, the angle bound gives uniform contraction for the
corresponding tree blocks (\Cref{cor:octahedral-block}).

\paragraph{Comparison with freely chosen random products.}
Amir, Blachar, Ghosh and Vishne~\cite{AmirMixability} study \emph{mixability}:
the existence of a finite random product that is exactly uniform, with its
steps and Bernoulli probabilities freely chosen. Every $S_n$ is mixable in
that sense; Angel--Holroyd~\cite{AngelHolroyd2018} construct exact shuffles
using lazy adjacent transpositions with suitable probabilities.
Our question instead fixes fair probabilities and a commutator
pattern, and asks for contraction uniformly over all valid configurations.
Among $S_n$ with $n\ge5$ and transposition leaves, only $S_5$ has such a
level. The distinction involves both the allowed products and the
quantifiers; the non-identity target is what makes the derived series
govern the threshold. \Cref{sec:related} also compares the construction
with word maps, conjugacy-class walks and expansion.

\paragraph{Origin.}
The motivation comes from Barrington's width-$5$ branching
programs~\cite{Bar89}: what is the distribution of their group product when
each expanded group instruction is independently replaced by the identity
with probability $1/2$? The companion paper~\cite{ZermenoS5} proves
target-independence and a quantitative mixing bound for the $S_5$ construction
on inconsistent reads. Here we establish the general structural law and
its quantitative consequences. For the conditioned law,
\Cref{thm:tindep,lem:conditioning} give
target-independence and an additive conditioning error. The fixed-group
mixing estimate uses the classical Diaconis--Shahshahani bound~\cite{DS81}.
Additional consequences include product-quotient obstructions, comparison
with a leaf Cayley walk, and determinant and congruence obstructions.

\paragraph{Reading guide.}
After the model and the conditioned branching-program law in \Cref{sec:prelim},
readers interested in the structural results can proceed to
\Cref{sec:criterion,sec:notfolklore} for \Cref{ithm:A,ithm:B}, and to
\Cref{subsec:rigid} for \Cref{ithm:C}.
\Cref{sec:layers,sec:synthesis,sec:rankone} give obstruction mechanisms,
examples and the classification of \Cref{ithm:D}.
The quantitative consequences and arithmetic application are developed in
\Cref{sec:quantitative-families,sec:octahedral-pair}.
\Cref{sec:computation} distinguishes exact certificates from numerical
checks, and \Cref{sec:open} records the remaining problems.

%% ===================================================================
\section{Preliminaries}
\label{sec:prelim}
%% ===================================================================

Let $G$ be a finite group. A \emph{permutation branching program} of length $L$ over
$G$ is a sequence of instructions $((i_1,g_1^{(0)},g_1^{(1)}),\dots)$; on input
$x\in\{0,1\}^\ell$ it outputs $\prod_{j=1}^L g_j^{(x_{i_j})}$. For a point function the
program is built from a fixed \emph{template} (a read sequence and a fixed
\emph{leaf element} $\sigma_j\in G$ per leaf): leaf $j$ outputs $\sigma_j$ when its
literal matches $T$ and the identity $e$ otherwise; the binary commutator tree of
Barrington's construction then combines the leaves.

\paragraph{Output laws.}
For a labelled tree $\tau$, drawing all expanded reads independently and
uniformly defines its \emph{unconditional} output law $\mu_\tau$.
For the associated point-function program we also write $P_T=\mu_\tau$:
the fixed template is implicit, and the unconditional law does not depend
on which branch matches $T$. Conditioning on inconsistent or consistent
paths gives $D_T$ or $C_T$, respectively.

\paragraph{Notation.}
Representations are unitary, with carrier space $V_\rho$ and dimension
$d_\rho$. We write $\widehat G$ for the irreducible representations and
$\mathbf1$ for the trivial one. For a probability measure $\mu$, put
$\widehat\mu(\rho)=\sum_g\mu(g)\rho(g)$ and
$M_\tau(\rho)=\widehat{\mu_\tau}(\rho)$.
The notation $\norm v$ denotes the Hilbert-space norm and $\norm A$ the
operator norm; $\frob A$ denotes the Frobenius norm.
We use $V_\rho^H$ for the $H$-fixed subspace and
$\Fix_\rho(s)=V_\rho^{\langle s\rangle}$ for a single element.
When $\rho$ is fixed, we may write $V=V_\rho$ and $V^H=V_\rho^H$.
For $w\in G$, its conjugacy class is denoted by $w^G$.

\paragraph{The independent-read commutator-tree model.}
The model has two structural features, recorded as \Cref{lem:factor,lem:tree}: the
$L$ bits are read \emph{independently}, and Barrington's program combines leaves through a
balanced \emph{commutator tree}. A level-$k$ target has $2^k$ configurable leaves; the expanded
independent-read program has $4^k$ leaf occurrences because inverse copies are read independently
with the inverse distribution. Write $\alpha(\tau)$ for the deterministic
commutator target of $\tau$ and $H(\tau)$ for the subgroup generated by its
leaf labels. The configuration is \emph{valid} when $\alpha(\tau)\ne e$.
For a conjugation-closed leaf set $\mathcal C$, a subgroup is $\mathcal C$-generated if it is
generated by its intersection with $\mathcal C$; the uniform strict-contraction level is the first
level at which every valid configurable $\mathcal C$-tree contracts in all non-trivial irreducible
representations.

For a transitive action of $G$ on $\Omega$ we write
$V=\C[\Omega]_0=\{f:\Omega\to\C\mid\sum_\omega f(\omega)=0\}$ for the \emph{augmentation
representation}. Its fixed spaces detect the orbits of the leaf subgroup,
as described in \Cref{prop:std}.

\begin{lemma}[Independent-read factorisation]
\label{lem:factor}
The operator $M_\tau(\rho)$ factors over the expanded leaf occurrences.
A point-function leaf with equiprobable branches $\{e,\sigma\}$ contributes the
\emph{leaf operator}
\[
P_\sigma(\rho)=\tfrac12\bigl(I+\rho(\sigma)\bigr),
\]
and for an involution leaf ($\sigma^2=e$) this is the \emph{orthogonal projection} onto
$\Fix_\rho(\sigma)$.
\end{lemma}
\begin{proof}
Independence makes the expectation defining $M_\tau(\rho)$ factor over the disjoint reads;
a leaf reading one bit with branches $e,\sigma$ of probability $\tfrac12$ each contributes
$\tfrac12(\rho(e)+\rho(\sigma))=\tfrac12(I+\rho(\sigma))$. If $\sigma^2=e$ then $\rho(\sigma)$ is a
self-adjoint unitary involution, so $P_\sigma$ is self-adjoint and
$P_\sigma^2=\tfrac14(I+\rho(\sigma))^2=\tfrac14(2I+2\rho(\sigma))=P_\sigma$.
\end{proof}

\begin{lemma}[Recursive commutator-tree operator]
\label{lem:tree}
Let $\tau=[\tau_1,\tau_2]$ be a node with its two child trees. Fix $\rho$
and put $X=M_{\tau_1}(\rho)$ and $Y=M_{\tau_2}(\rho)$. Then
\[
 M_\tau(\rho)=XYX^\dagger Y^\dagger.
\]
For a single leaf $s$, the operator is $P_s=P_s(\rho)$. Thus a node
with leaves $s,t$ has operator $P_sP_tP_s^\dagger P_t^\dagger$, equal
to $(P_sP_t)^2$ when both leaves are involutions.
Expanding the recursion gives a product of leaf averages and their
adjoints; involution leaf averages are orthogonal projections.
\end{lemma}
\begin{proof}
The four expanded sub-programs read disjoint positions. Their outputs
$a,b,a',b'$ are independent, with $a',b'$ distributed as the inverse
words $a^{-1},b^{-1}$. Unitarity gives
$\mathbb E\rho(a')=X^\dagger$ and $\mathbb E\rho(b')=Y^\dagger$.
Independence therefore yields
\[
 \mathbb E[\rho(a)\rho(b)\rho(a')\rho(b')]=XYX^\dagger Y^\dagger.
\]
For involution leaves, $P_s^\dagger=P_s$ and $P_t^\dagger=P_t$
by \Cref{lem:factor}.
\end{proof}

\begin{remark}[Independent occurrences and word evaluation]
\label{rem:notxy}
The two distributions
\[
\underbrace{XYX^{-1}Y^{-1}}_{\text{same }X,Y\ \text{(correlated)}}
\qquad\text{and}\qquad
\underbrace{XY\,\tilde X^{-1}\tilde Y^{-1}}_{\tilde X,\tilde Y\ \text{independent copies of }X,Y}
\]
encode different dependence between repeated occurrences. Ordinary word
evaluation uses the same value of $X$ in $X$ and $X^{-1}$. The
independent-read model uses four independent blocks, so its Fourier
transform factors as in \Cref{lem:tree}. This factorisation gives the
fixed-vector criterion of \Cref{lem:geom}.
\end{remark}

\paragraph{Mixing bound.}
The Diaconis--Shahshahani lemma~\cite{DS81} gives
\[
4\,\dtv(\mu,U_G)^2\le\sum_{\rho\neq\mathbf 1}d_\rho\frob{\widehat\mu(\rho)}^2,
\]
the sum over non-trivial irreducibles. Since $\frob{M}^2\le d_\rho\opn{M}^2$,
controlling $\opn{M_\tau(\rho)}$ for every non-trivial $\rho$ controls
$\dtv(\mu_\tau,U_G)$. Every total-variation upper bound below may also
be capped at $1$.

\subsection{The conditioned branching-program law}
\label{sec:tindep}
\label{sec:conditioning}

The structural and spectral results below concern the unconditional
law $\mu_\tau$. The following reductions from
\cite[Theorem~3 and Lemma~1]{ZermenoS5} transfer its mixing estimates
to conditioned branching-program laws.

\begin{theorem}[$T$-independence]
\label{thm:tindep}
Let $\BP_T,\BP_{T'}$ be point-function programs built from the same template (same read
sequence, same leaf element $\sigma_j$ at each leaf), differing only in which branch is
the matching one at the literals where $T_i\neq T'_i$. Then $D_T=D_{T'}$ over any group
$G$.
\end{theorem}

\begin{proof}
Let $S=\{i:T_i\neq T'_i\}$ and define $F_S:\{0,1\}^L\to\{0,1\}^L$ flipping $b_j$ exactly
when the literal read at step $j$ lies in $S$. Then (i) $F_S$ is an involution; (ii) it
flips all readings of each $S$-variable simultaneously, hence preserves the
consistency/inconsistency of each variable individually; and (iii) at every step the
swap of the matching branch ($T\!\leftrightarrow\!T'$) and the bit flip cancel, so the
output is unchanged: $\varphi_{T'}(b)=\varphi_T(F_S(b))$. As $F_S$ is a measure-preserving
involution of the inconsistent paths, change of variables gives $D_{T'}=D_T$.
\end{proof}

The same bijection proves $T$-independence over every finite group and
extends to programs with a fixed wildcard pattern.

\begin{lemma}[Conditioning]
\label{lem:conditioning}
Assume $0<p=\Pr[\text{consistent}]<1$, so both conditional laws are defined. Then $P_T=(1-p)D_T+p\,C_T$, hence
\[
\dtv(D_T,U_G)\;\le\;\dtv(P_T,U_G)+p\cdot\dtv(D_T,C_T)\;\le\;\dtv(P_T,U_G)+p .
\]
If every variable is read at least once then the consistent paths are in bijection with
$\{0,1\}^\ell$, so $p=2^\ell/2^L$.
\end{lemma}

\begin{proof}
The first identity is the law of total probability over the consistency event; the
triangle inequality and $\dtv\le1$ give the bound. The count of consistent paths is the
number of per-variable constant assignments.
\end{proof}

Any bound on $\dtv(P_T,U_G)$ therefore gives a bound for $D_T$ with additive
loss $p=2^{\ell-L}$. For $L=\ell^2$, this loss is
$2^{-\Theta(\ell^2)}$. We now estimate the unconditional mixing term.

%% ===================================================================
\section{The mixing criterion}
\label{sec:criterion}
%% ===================================================================

\subsection{When does an operator fail to contract?}

\begin{lemma}[Geometric lemma]
\label{lem:geom}
Fix a configuration $\tau$ with leaf set $S\subseteq G$ and let
$H=\langle S\rangle$. For any unitary representation $\rho$ on $V$,
\[
\opn{M_\tau(\rho)}=1\quad\Longleftrightarrow\quad V^H\neq0.
\]
\end{lemma}

\begin{proof}
Each leaf operator $P_\sigma(\rho)=\tfrac12(I+\rho(\sigma))$ is a contraction: $\rho(\sigma)$
is unitary with eigenvalues $e^{i\theta}$, so $P_\sigma(\rho)$ has eigenvalues
$\tfrac12(1+e^{i\theta})$ of modulus $\abs{\cos(\theta/2)}\le1$, equal to $1$ only on the
$+1$-eigenspace of $\rho(\sigma)$, i.e.\ on $\Fix_\rho(\sigma)$. By \Cref{lem:tree},
$M_\tau(\rho)$ is a product $M_1\cdots M_N$ in which each $M_t$ is a leaf operator
$P_\sigma(\rho)$ or its adjoint $P_\sigma(\rho)^{\dagger}=P_{\sigma^{-1}}(\rho)$ (some
$\sigma\in S$); both are normal contractions whose norm-$1$ eigenspace is exactly
$\Fix_\rho(\sigma)=\Fix_\rho(\sigma^{-1})$. If $\norm{M_1\cdots M_N v}=\norm{v}$ for a unit
vector $v$, put $v_N=v$ and $v_{t-1}=M_t v_t$. Since every $M_t$ is a contraction and
$\norm{v_0}=\norm{v_N}$, all inequalities $\norm{v_{t-1}}\le\norm{v_t}$ are equalities.
Equality for $M_t$ forces $v_t\in\Fix_\rho(\sigma_t)$, so $M_t v_t=v_t$ and
$v_{t-1}=v_t$. Inducting gives $v_t=v$ for all $t$, hence $v$ is fixed by every leaf, i.e.\
lies in $\bigcap_{\sigma\in S}\Fix_\rho(\sigma)=V^H$. Conversely any
$v\in V^H$ is fixed by $M_\tau(\rho)$. As $\opn{M_\tau(\rho)}\le1$ always, the claim
follows.
\end{proof}

\begin{theorem}[Fourier fixed-vector criterion (Theorem~B)]
\label{thm:criterion}
The Fourier operator of a configuration $\tau$ contracts strictly in \emph{every}
non-trivial irreducible representation \emph{if and only if its leaves generate $G$}:
\[
\opn{M_\tau(\rho)}<1\ \ \forall\,\rho\neq\mathbf1
\quad\Longleftrightarrow\quad
H(\tau)=G .
\]
\end{theorem}

\begin{proof}
Put $H=H(\tau)$. By \Cref{lem:geom} the left side fails for some non-trivial $\rho$ iff $V_\rho^H\neq0$
for some non-trivial $\rho$, i.e.\ iff $\Ind_H^G\mathbf1=\C[G/H]$ contains a non-trivial
irreducible. Now $\C[G/H]=\bigoplus_\rho (\dim V_\rho^H)\,\rho$ has dimension $[G:H]$ and
contains the trivial representation exactly once; it contains a non-trivial irreducible
iff $[G:H]>1$, i.e.\ iff $H\neq G$. Hence all non-trivial $\rho$ contract iff $H=G$.
\end{proof}

\begin{remark}
Combined with submultiplicativity
$\opn{ABA^\dagger B^\dagger}\le\opn A^2\opn B^2\le\max\{\opn A,\opn B\}^4$ and
the Diaconis--Shahshahani bound, \Cref{thm:criterion} says: in a fixed finite instance,
once every sub-block at some level $k_0$ has leaves generating $G$, every non-trivial
Fourier block is strictly contractive; taking the maximum over the finitely many
sub-blocks and irreducibles gives a gap $<1$, and further levels then decay. For a
family of programs, however, a quantitative total-variation bound also needs an explicit
lower bound on this gap and on the determinant/congruence layers.
\end{remark}

\subsection{The augmentation representation is only one test}

\begin{proposition}[Augmentation-representation specialisation]
\label{prop:std}
Assume $G$ acts transitively on $\Omega$. For a configuration $\tau$ and the
augmentation representation $V=\C[\Omega]_0$, $\opn{M_\tau(V)}=1$ iff
$H=H(\tau)$ is \emph{intransitive} on $\Omega$ (has $\ge2$ orbits). Consequently the
transitivity condition is \emph{necessary but not sufficient} for strict contraction
in every non-trivial irreducible: a proper
\emph{transitive} leaf subgroup contracts in $V$ yet, by \Cref{thm:criterion}, fails to
contract in some other irreducible.
\end{proposition}

\begin{proof}
The space
\[
V^{H}=\{f\in\C[\Omega]_0: f\text{ is constant on }H\text{-orbits}\}
\]
has dimension (number of $H$-orbits)$-1$, nonzero iff $H$ is intransitive; apply
\Cref{lem:geom}.
Intransitive $\Rightarrow H\neq G$, but the converse fails, and any proper transitive
$H$ witnesses insufficiency.
\end{proof}

For $S_5$ acting on the $6$ points of $\PGL_2(\F_5)$'s
projective line, a scan of the $27420$ reduced valid level-$2$ all-involution configurations
(\Cref{sec:computation}) finds the augmentation representation flags \emph{none}, yet $4695$
have leaves generating a proper transitive subgroup ($3705$ generating $A_5$, $990$ an
$S_4$-type subgroup of order $24$); these are invisible to $V$ but caught by the sign
representation (for the $A_5$ ones) and others. Over $S_5$ in its natural $5$-point
action the two tests happen to coincide because the leaves there are transpositions,
whose generated subgroups are Young subgroups; every proper such subgroup is intransitive.

\subsection{From the criterion to a solvability dichotomy}

\begin{lemma}[Validity is one-directional]
\label{lem:valid}
The level-$k$ commutator target of a configuration lies in $H^{(k)}$, the $k$-th derived
subgroup of $H=\langle S\rangle$ (the leaf subgroup). In particular, if $H^{(k)}=\{e\}$ the
target is trivial and the configuration is \emph{invalid}.
\end{lemma}
\begin{proof}
Induct on $k$. A level-$0$ target is a single leaf $\sigma\in S\subseteq H=H^{(0)}$. A
level-$(k{+}1)$ target is a commutator $[A,B]$ of two level-$k$ targets $A,B\in H^{(k)}$
(inductive hypothesis); hence $[A,B]\in[H^{(k)},H^{(k)}]=H^{(k+1)}$. The converse fails:
$H^{(k)}\neq\{e\}$ does not force a specific target to be non-trivial.
\end{proof}

By \Cref{thm:criterion} a contraction failure (in some irreducible) at level $k$ requires
a \emph{proper} leaf subgroup $H\neq G$ with $H^{(k)}\neq\{e\}$. To turn this into a
dichotomy we fix a \emph{conjugation-closed} leaf set $\mathcal{C}$ (a union of
conjugacy classes of $G$); all our examples---transpositions in $S_n$, involutions in
$A_5$ or $\PGL_2$---are of this form. The key fact is that, ranging over such a set, the
realisable targets generate the terms of the derived series.

\begin{lemma}[Targets generate the derived series; cf.~\cite{FernandezPintonello2023}]
\label{lem:central}
Let $\mathcal{C}\subseteq G$ be closed under conjugation by $H:=\langle\mathcal{C}\rangle$.
Define the realisable level-$k$ targets by $T_0=\mathcal{C}$ and
$T_{k+1}=\{[a,b]:a,b\in T_k\}$, where $a,b$ range independently (the left and right
sub-blocks of a balanced commutator tree read disjoint inputs and may be configured
independently). Then $\langle T_k\rangle=H^{(k)}$ for every $k\ge0$.
\end{lemma}
\begin{proof}
Induct on $k$; the base case is $\langle\mathcal{C}\rangle=H=H^{(0)}$. Conjugation by
$g\in H$ sends $[a,b]\mapsto[a^g,b^g]$ with $a^g,b^g\in T_k$, since $\mathcal{C}$---and
hence each $T_k$---is closed under $H$-conjugation; thus $T_{k+1}$ is
$H$-conjugation-closed and $\langle T_{k+1}\rangle$ is normal in $H$. Write
$M=H^{(k)}=\langle T_k\rangle$. Each $[a,b]$ ($a,b\in T_k\subseteq M$) lies in
$[M,M]=H^{(k+1)}$, so $\langle T_{k+1}\rangle\subseteq H^{(k+1)}$. Conversely, $[M,M]$ is
the normal closure in $M$ of $\{[s,t]:s,t\in T_k\}$ (in $M/N$, with $N$ that normal
closure, the images of the generating set $T_k$ commute, so $M/N$ is abelian); as
$\langle T_{k+1}\rangle$ is normal in $M$ and contains those commutators, it contains
$[M,M]=H^{(k+1)}$. Hence equality.
\end{proof}

\begin{remark}[Configurable leaf sets and rigid templates]
\label{rem:configurable}
In the configurable model, the two sub-blocks range independently over
the conjugation-closed leaf set. Their targets generate $H^{(k)}$ by
\Cref{lem:central}, giving the exact threshold in \Cref{thm:dichotomy}.
A rigid template specifies a leaf at each position; its realised targets
form a restricted subset of the configurable targets. The same value $D$
then gives an upper bound. \Cref{subsec:rigid} studies attainable rigid
levels and persistence of valid binary continuations.
\end{remark}

\begin{definition}[$\mathcal{C}$-generated subgroup]
\label{def:cgen}
A subgroup $H\le G$ is \emph{$\mathcal{C}$-generated} if $H=\langle H\cap\mathcal{C}\rangle$.
Every leaf subgroup $L=\langle S\rangle$ with $S\subseteq\mathcal{C}$ is
$\mathcal{C}$-generated; and, $\mathcal{C}$ being conjugation-closed in $G$,
$L\cap\mathcal{C}$ is closed under conjugation by $L$.
\end{definition}

\begin{theorem}[Solvability dichotomy and exact level (Theorem~A)]
\label{thm:dichotomy}
Let $G$ be a finite \emph{non-solvable} group and $\mathcal{C}\subseteq G$ a
conjugation-closed leaf set with $\langle\mathcal{C}\rangle=G$, in the \emph{configurable} model
where the two sub-blocks of each node may be configured independently (\Cref{rem:configurable}; for
a single rigid template the level $D$ below is only an upper bound). {\normalfont(A solvable
ambient is degenerate: no valid configuration of large level exists, its target lying in
$G^{(d)}=\{e\}$.)}
\begin{enumerate}[nosep,leftmargin=1.6em]
\item\textbf{(Obstruction.)} If some proper $\mathcal{C}$-generated subgroup is
non-solvable, then for every level $k$ a valid configuration $\tau$ has $\opn{M_\tau(\rho)}=1$
for some non-trivial $\rho$: the leaf set has \emph{no uniform strict-contraction level}.
\item\textbf{(Sharp level.)} If every proper $\mathcal{C}$-generated subgroup is solvable,
put $D=\max\{\operatorname{dl}(H):H<G,\ H\text{ is }\mathcal{C}\text{-generated}\}$. Then
every valid configuration of level $\ge D$ has leaf subgroup $G$ and contracts strictly in
all non-trivial irreducibles, while some valid configuration of level $D-1$ does not. The
level of strict contraction is exactly $D$.
\end{enumerate}
\end{theorem}
\begin{proof}
\emph{(1)} Let $H<G$ be proper, $\mathcal{C}$-generated and non-solvable; its derived
series stabilises at a non-trivial perfect core, so $H^{(k)}\neq\{e\}$ for all $k$. By
\Cref{lem:central} inside $H$ (with the $H$-conjugation-closed set $H\cap\mathcal{C}$,
which generates $H$ by \Cref{def:cgen}), the level-$k$ targets with leaves in
$H\cap\mathcal{C}$ generate $H^{(k)}\neq\{e\}$, so a valid level-$k$ configuration $\tau$ with
leaves in $H$ exists. Its leaf subgroup is $\le H<G$, so by \Cref{thm:criterion} some
non-trivial $\rho$ has $V_\rho^{H(\tau)}\neq0$, whence
$\opn{M_\tau(\rho)}=1$ (\Cref{lem:geom}). \emph{(2)} If a valid level-$k$ configuration
($k\ge D$) had leaf subgroup $L\neq G$, then $L$ is proper, $\mathcal{C}$-generated and,
by hypothesis, solvable with $\operatorname{dl}(L)\le D\le k$, so its target lies in
$L^{(k)}\subseteq L^{(D)}=\{e\}$ (\Cref{lem:valid})---contradicting validity. Thus $L=G$
and it contracts strictly (\Cref{thm:criterion}). For exactness, take proper
$\mathcal{C}$-generated $H$ with $\operatorname{dl}(H)=D$; by \Cref{lem:central} inside $H$
the level-$(D-1)$ targets generate $H^{(D-1)}\neq\{e\}$, giving a valid level-$(D-1)$
configuration with leaf subgroup $\le H<G$ that fails to contract (\Cref{lem:geom},
\Cref{thm:criterion}).
\end{proof}

The contraction level is therefore exactly the largest derived length among proper
$\mathcal{C}$-generated subgroups, and yields a quantitative bound for each fixed group.
The value $D$ is also the first level at which uniform contraction holds: if a proper
$\mathcal C$-generated $H$ has $\operatorname{dl}(H)=D$, then $H^{(k)}\neq\{e\}$ for every
$k<D$, and \Cref{lem:central} inside $H$ supplies a valid configuration with leaf subgroup
contained in $H$ at each of these levels, so strict contraction fails at every level below $D$.

\begin{proposition}[Quantitative bound, fixed group]
\label{prop:quant}
In case (2) of \Cref{thm:dichotomy}, let $\lambda=\lambda_{G,\mathcal C}<1$
be the maximum of $\opn{M_\tau(\rho)}$ over non-trivial irreducibles
$\rho$ and valid level-$D$ configurations $\tau$.
For every valid configuration $\tau$ of level $m\ge D$,
\[
 \dtv(\mu_\tau,U_G)\le
 \min\Bigl\{1,\tfrac12\sqrt{|G|-1}\,\lambda^{4^{m-D}}\Bigr\}.
\]
For the associated conditioned law $D_T$, add $p$ from
\Cref{lem:conditioning} to this bound.
\end{proposition}
\begin{proof}
Every sub-block of a valid configuration is valid: were a sub-block target $e$, each
ancestor commutator $[e,\cdot]=e$ would vanish, forcing the root target to be $e$. Hence
each level-$D$ sub-block has leaf subgroup $G$ and operator norm $\le\lambda$
(case~(2)), and submultiplicativity gives
$\opn{M_\tau(\rho)}\le\lambda^{4^{m-D}}$.
With $\frob{M}^2\le d_\rho\opn{M}^2$ and $\sum_\rho d_\rho^2=|G|$, Diaconis--Shahshahani
gives $4\dtv(\mu_\tau,U_G)^2\le(|G|-1)\lambda^{2\cdot4^{m-D}}$; the $D_T$ bound is
\Cref{lem:conditioning}. For $G=S_5$ the prefactor $\tfrac12\sqrt{119}\approx5.45$
recovers the companion paper's bound. \Cref{sec:quantitative-families}
controls the gap as the group varies and derives uniform family rates.
\end{proof}

\subsection{Functoriality and localisation of obstructions}

The dichotomy is monotone under quotients: any obstruction visible in a quotient is already an
obstruction upstairs. Thus, to mix, $G$ must mix through every non-solvable quotient; solvable
quotients are the degenerate case already noted in \Cref{thm:dichotomy}.

\begin{proposition}[Quotients]
\label{prop:quotient}
Let $N\trianglelefteq G$ and $\bar{\mathcal C}=\pi(\mathcal C)\subseteq G/N$ its image under
$\pi:G\to G/N$ (again conjugation-closed). If $(G,\mathcal C)$ is in the mixing case of
\Cref{thm:dichotomy}, then every proper $\bar{\mathcal C}$-generated subgroup of $G/N$ is
solvable. In particular, if $G/N$ is non-solvable, then $(G/N,\bar{\mathcal C})$ is also in the
mixing case, and the level cannot rise:
\[
D(G/N,\bar{\mathcal C})\le D(G,\mathcal C).
\]
Equivalently, an obstruction in a quotient lifts to one in $G$.
\end{proposition}
\begin{proof}
Let $\bar K<G/N$ be proper and $\bar{\mathcal C}$-generated, and set
$L=\langle\mathcal C\cap\pi^{-1}(\bar K)\rangle$. Then $L$ is $\mathcal C$-generated, and since
$\pi\bigl(\mathcal C\cap\pi^{-1}(\bar K)\bigr)=\bar{\mathcal C}\cap\bar K$ we have
$\pi(L)=\langle\bar{\mathcal C}\cap\bar K\rangle=\bar K$. Hence $L<G$ is proper and
$\operatorname{dl}(L)\ge\operatorname{dl}(\pi(L))=\operatorname{dl}(\bar K)$. So a non-solvable
$\bar K$ forces a non-solvable $L$ (the obstruction lifts); and if every proper $\mathcal
C$-generated subgroup of $G$ is solvable, so is $L$, whence $\bar K$ is solvable with
$\operatorname{dl}(\bar K)\le\operatorname{dl}(L)\le D(G,\mathcal C)$. Maximising over $\bar K$
gives the level bound when $G/N$ is non-solvable.
\end{proof}

\begin{corollary}[An obstruction from product quotients]
\label{cor:product-obstruction}
Suppose a finite group $G$ has a quotient $A\times B$, with $A$ non-solvable
and $B$ non-cyclic. Every conjugation-closed generating leaf set of $G$
then has a permanent obstruction in the configurable model.
\end{corollary}
\begin{proof}
First suppose $G=A\times B$ and let $\mathcal C$ generate $G$.
The first projection of $\mathcal C$ generates $A$, so some $(a,b)\in\mathcal C$
has $a\notin O_\infty(A)$. The normal closure of $a$ in $A$ is non-solvable:
otherwise it would be contained in the solvable radical.
Conjugation-closure implies
\[
S=\{(xax^{-1},b):x\in A\}\subseteq\mathcal C.
\]
The group $H=\langle S\rangle$ projects onto that non-solvable normal closure,
while $H\le A\times\langle b\rangle<A\times B$ since $B$ is non-cyclic.
Thus $H$ is a proper non-solvable $\mathcal C$-generated subgroup.
For a quotient of $G$, apply this argument to the image leaf set and
then lift the obstruction by \Cref{prop:quotient}.
\end{proof}

The non-cyclic hypothesis cannot simply be omitted. For example, let
$G=A_5\times C_2$ and take
$\mathcal C=\{(t,z):t\in A_5\text{ an involution}\}$, where $z$ generates $C_2$.
This class generates $G$: a subgroup projecting onto $A_5$ with trivial
kernel would be the graph of a homomorphism $A_5\to C_2$, which cannot
contain such leaves. Similarly, a proper $\mathcal C$-generated subgroup
cannot project onto all of $A_5$ and is therefore contained in $K\times C_2$
for a proper involution-generated $K<A_5$. These $K$ have derived length
at most $2$ by \Cref{cor:A5}; the bound is attained by lifting the
involutions of an $S_3$ along its sign character. Hence this class has
exact level $2$. In contrast, $A_5\times(C_2\times C_2)$ admits no mixing
conjugation-closed generating class.


\begin{proposition}[Perfect-core localisation]
\label{prop:perfectcore}
For a finite group write $H^{(\infty)}=\bigcap_iH^{(i)}$ for the terminal (perfect) term of the
derived series, so $H$ is solvable iff $H^{(\infty)}=1$. Let $N\trianglelefteq G$,
$\pi\colon G\to G/N$, and let $H<G$ be a proper non-solvable $\mathcal C$-generated subgroup
(an obstruction in \Cref{thm:dichotomy}) with $HN\neq G$. Then $\pi(H)$ is a proper
$\bar{\mathcal C}$-generated subgroup, and exactly one holds:
\begin{enumerate}[nosep,leftmargin=1.6em]
\item $\pi(H)$ is non-solvable --- the obstruction \emph{descends}, $\pi(H)$ being an obstruction
for $(G/N,\bar{\mathcal C})$; or
\item $\pi(H)$ is solvable --- the perfect core falls into the kernel, $H^{(\infty)}\le H\cap N$.
\end{enumerate}
\end{proposition}
\begin{proof}
$HN\neq G$ gives $\pi(H)=HN/N<G/N$ proper, and $H=\langle H\cap\mathcal C\rangle$ gives
$\pi(H)=\langle\pi(H\cap\mathcal C)\rangle$, hence $\bar{\mathcal C}$-generated. Since
$\pi(H^{(i)})=\pi(H)^{(i)}$ and a finite derived series stabilises,
$\pi(H^{(\infty)})=\pi(H)^{(\infty)}$. If $\pi(H)$ is solvable then $\pi(H)^{(\infty)}=1$, so
$H^{(\infty)}\le\ker\pi\cap H=H\cap N$; otherwise $\pi(H)$ is non-solvable. (Should a leaf lie in
$N$, $e\in\bar{\mathcal C}$; discarding it changes no generated subgroup.)
\end{proof}

\begin{remark}[Schreier generators for kernel obstructions]
\label{rem:schreier}
In case~(ii), the kernel $K=H\cap N$ may contain no leaves. For example,
for $G=S_5\times C_2$, $N=A_5\times1$, leaves the transpositions of $S_5$ and a
generator of $C_2$, the obstruction $H=S_5\times1$ has $H^{(\infty)}=A_5\times1\le N$ while
$\mathcal C\cap N=\varnothing$. Reidemeister--Schreier supplies generators: for
$H=\langle S\rangle$ and $K=H\cap N$, choosing representatives $t\in H$ for the cosets
$K\backslash H\simeq\pi(H)$ gives generators
\[
t\,s\,\overline{ts}^{-1}\qquad (s\in S),
\]
where $\overline{ts}$ denotes the chosen representative of the coset $Kts$. If the representatives
are chosen by shortest words in the image generating set $\pi(S)$, these Schreier words have length
at most $2\,\mathrm{diam}(\pi(H),\pi(S))+1$ in the original leaves. The kernel inherits the
leaves' \emph{lifting errors}, providing a mechanism for
congruence-layer obstructions in $\GL_n(\Z/p^r)$ (\Cref{sec:layers}).
\end{remark}

\begin{proposition}[Clean reduction over a Frattini solvable kernel]
\label{prop:frattini}
Let $N\trianglelefteq G$ be solvable with $N\le\Phi(G)$, the Frattini subgroup. Then no proper
subgroup supplements $N$ ($HN=G\Rightarrow H=G$), and in the non-degenerate regime
($G$ non-solvable and $\langle\mathcal C\rangle=G$, hence $G/N$ non-solvable and
$\langle\bar{\mathcal C}\rangle=G/N$)
\[
\begin{aligned}
(G,\mathcal C)\ \text{mixes}
  &\iff(G/N,\bar{\mathcal C})\ \text{mixes},\\
\text{in the mixing case:}\quad D(G/N,\bar{\mathcal C})
  &\le D(G,\mathcal C)
  \le D(G/N,\bar{\mathcal C})+\operatorname{dl}(N).
\end{aligned}
\]
\end{proposition}
\begin{proof}
If $H<G$ had $HN=G$, then $H\le M$ for some maximal $M\supseteq\Phi(G)\supseteq N$, whence
$G=HN\le M<G$; so $HN=G$ forces $H=G$, and every proper $\mathcal C$-generated $H$ has $\pi(H)$
proper. If $H$ is a non-solvable obstruction then $H^{(\infty)}\neq1$, while $H^{(\infty)}\le N$
would make it a perfect subgroup of the solvable $N$ and force $H^{(\infty)}=1$; so case~(ii) of
\Cref{prop:perfectcore} is impossible and $\pi(H)$ is a non-solvable obstruction in $G/N$. With
\Cref{prop:quotient} this gives the equivalence. For the level, a proper solvable $\mathcal
C$-generated $H$ has $\operatorname{dl}(H)\le\operatorname{dl}(\pi(H))+\operatorname{dl}(H\cap
N)\le D(G/N,\bar{\mathcal C})+\operatorname{dl}(N)$ (with $\operatorname{dl}(1)=0$); the left
inequality is \Cref{prop:quotient}.
\end{proof}

\begin{remark}[Supplements to the kernel]
\label{rem:residue}
A proper supplement $H<G$ with $HN=G$ maps onto the whole quotient.
When $G/N$ is non-solvable, so is $H$, and detecting this obstruction
requires information about supplements. \Cref{cor:product-obstruction}
gives one such obstruction for product quotients.
\end{remark}

\subsection{Stability of the criterion}

The fixed-vector criterion has a quantitative form: a vector that loses
little norm is close to the fixed space of the leaf subgroup. The estimate
applies to leaves of arbitrary order.

\begin{proposition}[Stability / quantitative gap]
\label{prop:stability}
Let $M_j=(I+\rho(s_j))/2$, $T=M_N\cdots M_1$, and
$H=\langle s_1,\ldots,s_N\rangle\ne\{e\}$, where $\rho$ is unitary on $V$ and
the $s_j$ may have arbitrary orders. Let $L$ be the diameter of $H$ in
$\{s_1,\ldots,s_N\}^{\pm1}$. For a unit vector $v$, put
$E=1-\norm{Tv}^2$. Then
\[
 \norm{(I-M_j)v}\le\sqrt{NE},\qquad
 \operatorname{dist}(v,V^H)\le2L\sqrt{NE}.
\]
In particular, if $V^H=0$, then
\begin{equation}\label{eq:diameter-gap}
 \opn T^2\le1-\frac1{4NL^2},\qquad
 \opn T\le1-\frac1{8NL^2}.
\end{equation}
\end{proposition}
\begin{proof}
For a unitary $U$ and $M=(I+U)/2$, the parallelogram identity gives
$\norm{Mx}^2+\norm{(I-M)x}^2=\norm x^2$, and $\opn{I-M}\le1$.
Set $w_0=v$, $w_j=M_jw_{j-1}$ and $e_j=w_{j-1}-w_j$. Thus
$\sum_j\norm{e_j}^2=E$ and
\[
 \norm{(I-M_j)v}\le\norm{e_j}+\norm{v-w_{j-1}}
 \le\sum_{i\le j}\norm{e_i}\le\sqrt{NE}.
\]
Consequently $\norm{(I-\rho(s_j))v}\le2\sqrt{NE}$; the same holds
for $s_j^{-1}$. Telescoping a word of length at most $L$ gives
$\norm{(I-\rho(h))v}\le2L\sqrt{NE}$ for every $h\in H$.
Averaging over $H$ proves the distance bound. If $V^H=0$, this yields
$1\le2L\sqrt{NE}$ for every unit $v$, proving the first inequality in
\eqref{eq:diameter-gap}; use $\sqrt{1-x}\le1-x/2$ for the second.
\end{proof}

\begin{remark}
For generating leaves in a non-trivial irreducible, \Cref{prop:stability}
gives an explicit gap in terms of their Cayley diameter. With $N=4^d$
at level $d$, this estimate yields the family rates in
\Cref{sec:quantitative-families}.
\end{remark}

\begin{proposition}[Comparison with the leaf Cayley walk]
\label{prop:cayley-comparison}
Let $s_1,\ldots,s_N$ be involutions in a finite group $G$, listed with
multiplicity and $N\ge1$. On $\ell^2_0(G)$ let $R$ be the regular
representation and put
\[
P_i=\tfrac12(I+R(s_i)),\qquad T=P_N\cdots P_1,\qquad
A=\frac1N\sum_{i=1}^N P_i,\qquad \delta=1-\opn A.
\]
Then
\[
\frac{2\delta}{N+1}\ \le\ 1-\opn T^2\ \le\ \min\{1,4N\delta\}.
\]
Here $A$ is the positive semidefinite transition operator of the lazy
Cayley walk with the indicated multiplicities. The assertions include
$\delta=0$, which occurs exactly when the leaves generate a proper subgroup.
\end{proposition}
\begin{proof}
For a unit vector $v$ write $w_0=v$, $w_i=P_iw_{i-1}$ and
$e_i=(I-P_i)w_{i-1}$. Orthogonality gives
$\sum_i\norm{e_i}^2=1-\norm{Tv}^2$. Also
\[
\norm{(I-P_i)v}\le\sum_{j\le i}\norm{e_j},\qquad
\sum_i\norm{(I-P_i)v}^2
\le\frac{N(N+1)}2\sum_j\norm{e_j}^2.
\]
Since $I-A\ge\delta I$, the left sum is at least $N\delta$.
Choose $v$ attaining $\opn T$ to obtain the lower bound.
For the upper bound choose a unit eigenvector $v$ of $I-A$ with
eigenvalue $\delta$. Gao's projection union bound, in the form
\cite[Lemma 2]{AnshuAradVidick2016}, gives
\[
\norm{Tv}^2\ge1-4\sum_i\norm{(I-P_i)v}^2=1-4N\delta.
\]
The additional bound by $1$ follows from $\opn T^2\ge0$.
\end{proof}

\begin{corollary}[Bounded-depth uniformity]
\label{cor:bounded-depth-gap}
For a family of involutive independent-read blocks with uniformly bounded
expanded length $N$, their operator norms on $\ell^2_0(G)$ are uniformly
bounded below $1$ if and only if the associated lazy Cayley walks have
uniformly positive spectral gaps. In a level-$k$ block $N=4^k$;
the walk uses the expanded occurrences, equivalently the $2^k$ leaf
positions with their multiplicities. Thus the assertion applies to
uniformly bounded depths, and asserts no such equivalence with uniform
constants when depths are unbounded.
\end{corollary}
\begin{proof}
Both implications follow from \Cref{prop:cayley-comparison}; the regular
representation contains every non-trivial irreducible, so its norm on
$\ell^2_0(G)$ is the maximum of their norms.
\end{proof}


\subsection{Averaging leaves}

The fixed-vector criterion extends to general lazy averages. For a
probability distribution $\mu$ on $G$ with $e\in\operatorname{supp}\mu$,
put $P_\mu(\rho)=\mathbb E_{g\sim\mu}[\rho(g)]=\sum_g\mu(g)\rho(g)$.

\begin{proposition}[Averaging leaves]
\label{prop:general-leaf}
For lazy $\mu$, $\opn{P_\mu}\le1$, and $\norm{P_\mu v}=\norm v$ iff $v\in V_\rho^{H_\mu}$ with
$H_\mu=\langle\operatorname{supp}\mu\rangle$ (and then $P_\mu v=v$). Hence
\Cref{lem:geom,thm:criterion} hold verbatim with $H_\mu$ in place of $\langle\sigma\rangle$:
averaging leaves contract in every non-trivial irreducible iff
$\langle\bigcup_{\text{leaves}}H_\mu\rangle=G$. The point-function leaf
$\mu=\mathrm{unif}\{e,\sigma\}$ is the case $H_\mu=\langle\sigma\rangle$.
\end{proposition}
\begin{proof}
$\norm{\sum_g\mu(g)\rho(g)v}\le\sum_g\mu(g)\norm{\rho(g)v}=\norm v$, with equality iff the equal-norm
vectors $\rho(g)v$ ($g\in\operatorname{supp}\mu$) coincide; as $e\in\operatorname{supp}\mu$ they then
all equal $\rho(e)v=v$, i.e.\ $v\in V_\rho^{H_\mu}$, and $P_\mu v=\sum_g\mu(g)v=v$. The peeling of
\Cref{lem:geom} is unchanged.
\end{proof}

\begin{remark}
For these averaging operators, the generated supports determine the
common fixed space. The target and validity conditions in
\Cref{thm:dichotomy} apply when the underlying leaf labels range over a
configurable conjugation-closed class.
\end{remark}

\subsection{Subgroup averages and finite Hecke boundary operators}
\label{subsec:hecke-boundary}

Subgroup averages connect fixed spaces to the angles used in
\Cref{sec:octahedral-pair} through finite Hecke operators.
Fix a unitary representation $\rho$ on $V$. In this subsection we write
$P_H=P_H^\rho$ for the projection onto $V^H$.

\begin{lemma}[Subgroup averages are fixed-space projections]
\label{lem:subgroup-average}
Let $H\le G$ and set
\[
  e_H=\frac1{\abs H}\sum_{h\in H}h\in\C[G],
  \qquad
  P_H=\rho(e_H)=\frac1{\abs H}\sum_{h\in H}\rho(h).
\]
Then $P_H$ is the orthogonal projection onto $V^H$.
\end{lemma}

\begin{proof}
The average is self-adjoint and idempotent. Its image is $H$-fixed by
reindexing the sum, and it acts as the identity on every $H$-fixed vector.
\end{proof}

\begin{proposition}[Subgroup averaging and finite Hecke compression]
\label{prop:hecke-compression}
Let $H\le G$, let $g\in G$, and let $U_H$ be the uniform probability measure on $H$. Then
\[
  \widehat{U_H*\delta_g*U_H}(\rho)=P_H\rho(g)P_H .
\]
The measure $U_H*\delta_g*U_H$ depends only on the double coset $HgH$, and the corresponding
group-algebra element belongs to the finite Hecke corner $e_H\C[G]e_H$.
\end{proposition}

\begin{proof}
Convolution becomes multiplication under $\rho$. For $h_1,h_2\in H$,
$e_Hh_1gh_2e_H=e_Hge_H$, proving the double-coset assertion.
\end{proof}

\begin{corollary}[Double cosets as angles between fixed spaces]
\label{cor:double-coset-angle}
Let $K=gHg^{-1}$. Then
\[
  P_K=\rho(g)P_H\rho(g)^{-1},
  \qquad
  P_H\rho(g)P_H=P_HP_K\rho(g),
\]
and therefore
\[
  \opn{P_H\rho(g)P_H}=\opn{P_HP_K}.
\]
Thus the norm of the double-coset compression is exactly the cosine of the smallest principal
angle between the fixed spaces $V^H$ and $V^K$, with cosine defined as $0$ when either space is zero.
\end{corollary}

\begin{proof}
Conjugate the average over $H$, then use $P_K\rho(g)=\rho(g)P_H$.
Right multiplication by the unitary $\rho(g)$ preserves the norm.
\end{proof}

If $\langle H,K\rangle=G$ and $\rho$ is non-trivial irreducible, then
$V^H\cap V^K=V^G=0$. The angle in
\Cref{cor:double-coset-angle} is therefore the Friedrichs angle. This is
the angle controlling classical alternation of two orthogonal projections;
see~\cite{BadeaGrivauxMuller2010} for that formula and its distinction from
products involving more than two subspaces. The tree estimates here use
their stated product bounds, and \Cref{lem:octahedral-gluing} also allows
blocks that only approximate fixed-space projections.

\begin{lemma}[Local-to-boundary stability]
\label{lem:local-boundary-stability}
Let $H,K\le G$, let $g\in G$, and let $\nu_H,\nu_K$ be probability measures supported on $H,K$.
Put $X=\widehat{\nu_H}(\rho)$ and $Y=\widehat{\nu_K}(\rho)$. If
\[
  \opn{X-P_H}\le\varepsilon_H,
  \qquad
  \opn{Y-P_K}\le\varepsilon_K,
\]
then
\[
  \opn{X\rho(g)Y-P_H\rho(g)P_K}\le\varepsilon_H+\varepsilon_K .
\]
\end{lemma}

\begin{proof}
Insert $P_H\rho(g)Y$ between the two terms.
The triangle inequality and the fact that all other factors are contractions
give $\varepsilon_H+\varepsilon_K$.
\end{proof}

\begin{remark}
An internal contraction estimate controls the error in replacing a block
by a subgroup average; the angle then controls the interaction of two
blocks. \Cref{lem:octahedral-gluing,thm:octahedral-sharp} implement this
for the octahedral pairs with intersection of order six.
\end{remark}

\subsection{Effectivity}

\begin{proposition}[Decidability]
\label{prop:algorithm}
Given an effective finite group model in which subgroups can be enumerated (for example, a
multiplication table, or permutation groups via standard algorithms) and a conjugation-closed leaf
set $\mathcal C$, it is decidable whether $(G,\mathcal C)$ is in the mixing case of
\Cref{thm:dichotomy}; when it is, the exact level $D$ is computable, and otherwise a non-solvable
obstructing subgroup is exhibited (or the ambient/input is reported as degenerate, e.g.
$\langle\mathcal C\rangle\neq G$ or $G$ solvable).
\end{proposition}
\begin{proof}
First check whether $\langle\mathcal C\rangle=G$ and whether the ambient is non-solvable. Then
enumerate the subgroups $H<G$, retain exactly the $\mathcal C$-generated ones
($H=\langle H\cap\mathcal C\rangle$), and test each for solvability by its derived series. Taking
the maximum derived length among the solvable witnesses returns $D$; a non-solvable one is the
obstruction in \Cref{thm:dichotomy}.
\end{proof}

\begin{remark}
The cost of the procedure is governed by the number of subgroups enumerated.
\end{remark}

The worked instances---$S_n$, $A_5$, $\PGL_2(\F_q)$---are collected in
\Cref{sec:synthesis}; the matrix groups $\GL_n(\F_q)$ reduce to a subgroup question
(\Cref{sec:layers}).

\subsection{The analytic--algebraic correspondence}
\label{sec:notfolklore}

The expanded independent-read operator is a product of lazy averages. Its
norm-one condition depends on the subgroup generated by the leaves, as in
\Cref{lem:geom,thm:criterion}. The derived series enters through the additional
requirement that the folded commutator target be non-identity. The normal-subset
word-value lemma~\cite[Lemma 2.1(ii)]{FernandezPintonello2023} identifies the subgroup
generated by the realisable depth-$k$ targets with $H^{(k)}$.
Combining these two facts gives the exact obstruction and level of
\Cref{thm:dichotomy}.

The quantitative correspondence is recorded in
\Cref{prop:cayley-comparison}: for involution leaves it compares the block
gap with that of the associated lazy Cayley walk.
\Cref{tab:dictionary} collects the connections used throughout the paper.

\begin{table}[htbp]
\centering\small
\begin{tabular}{p{.30\textwidth}p{.62\textwidth}}
\toprule
Tree quantity & Correspondence \\
\midrule
Leaf average & A lazy group step; an orthogonal fixed-space projection for an involution. \\
Strict contraction & Generation by the leaves, when tested in every non-trivial irreducible (\Cref{ithm:B}). \\
First uniform level & Maximal proper leaf-subgroup derived length, in the unobstructed configurable case (\Cref{ithm:A}). \\
Uniform block gap & Expansion of the associated lazy Cayley walks, for involutions at bounded depth (\Cref{cor:bounded-depth-gap}). \\
Two subgroup averages & Fixed-space angle; for conjugate subgroups, the norm of a finite Hecke compression (\Cref{cor:double-coset-angle}). \\
\bottomrule
\end{tabular}
\caption{Connections used in the structural and quantitative arguments.}
\label{tab:dictionary}
\end{table}

%% ===================================================================
\section{The three layers and matrix obstructions}
\label{sec:layers}
%% ===================================================================

For $G=\GL_n(\F_q)$ the determinant $\det:G\twoheadrightarrow\F_q^\times$
is an abelian quotient with kernel $\SL_n$. Three sources of obstructions
are useful in these matrix families: subspace stabilisers, the determinant
quotient, and congruence kernels over $\Z_{p^k}$. We examine each through
the leaf-subgroup criterion.

\paragraph{(a) The shape, and matrix obstructions.}
The non-abelian content is \Cref{thm:criterion} and \Cref{thm:dichotomy}. For the relevant
matrix families, parabolics (subspace stabilisers) give an explicit family of
obstructions. The stabiliser of an $r$-subspace of $\F_q^n$ has Levi
factor $\GL_r\times\GL_{n-r}$, which is non-solvable as soon as one factor is. A leaf set whose
elements can remain inside such a parabolic may therefore contain a proper non-solvable
$\mathcal C$-generated subgroup.

\begin{proposition}[Parabolic obstruction in the matrix regime]
\label{prop:parabolic}
Let $P<\GL_n(\F_q)$ be a proper parabolic and let
$H_P=\langle P\cap\mathcal C\rangle$. If the image of $H_P$ in a Levi factor of $P$ is
non-solvable, then $H_P$ is a proper non-solvable $\mathcal C$-generated subgroup. Consequently
the leaf set has a permanent obstruction in the sense of \Cref{thm:dichotomy}(1). Thus
parabolics provide explicit structural obstructions whenever this Levi-image condition can be
verified. By contrast, $\GL_2(\F_q)$ has only the Borel as a proper parabolic, and the Borel is
solvable; $\GL_3(\F_2)$ has Levi $\GL_2(\F_2)\cong S_3$ (solvable), so this particular parabolic
obstruction does not appear in these small cases.
\end{proposition}

Parabolics furnish an explicit infinite family of obstructions; the remaining maximal subgroups
reduce to a Dickson/Aschbacher generation question (\Cref{sec:open}).

\begin{proposition}[Transvection leaves]
\label{prop:transvection-levels}
Let $n\ge3$, let $G=\SL_n(q)$, and let $\mathcal T$ be its class of
transvections. The exact configurable contraction level is $3$ for
$(n,q)=(3,2)$ and $4$ for $(n,q)=(3,3)$. In every other case there is
a permanent obstruction.
\end{proposition}
\begin{proof}
Transvections generate $\SL_n(q)$. For $q\ge4$ the upper-left
$\SL_2(q)$ is a proper non-solvable $\mathcal T$-generated subgroup.
For $q\in\{2,3\}$ and $n\ge4$, use the upper-left $\SL_3(q)$ instead.
Apply \Cref{thm:dichotomy}(1). In $\SL_3(2)$ transvections are precisely
the involutions, and \Cref{rem:gl32} gives level $3$.

For $\SL_3(3)=\PSL_3(3)$ use the maximal-subgroup list in
\Cref{prop:psl33-level}. In a point parabolic $P=3^2{:}\GL_2(3)$,
a transvection acts trivially on the invariant line, so its Levi image
lies in $\SL_2(3)$. The radical $V\cong\F_3^2$ and the elementary
matrices in this Levi factor are transvections; hence
$\langle P\cap\mathcal T\rangle=J=V\rtimes\SL_2(3)$.
Since $\SL_2(3)'=Q_8$, $Q_8'=\langle-I\rangle$, and $-I$ acts as
negation on $V$, the derived series is
\[
 J\ >\ V\rtimes Q_8\ >\ V\rtimes\langle-I\rangle\ >\ V\ >\ 1,
\]
with orders $216,72,18,9,1$. The dual parabolic has the same property.
The other maximals, $S_4$ and $13{:}3$, have derived lengths $3$ and $2$.
Every proper $\mathcal T$-generated subgroup therefore has derived
length at most $4$, and $J$ attains $4$. Apply \Cref{thm:dichotomy}(2).
\end{proof}

\paragraph{(b) The determinant layer.}
Apart from the small exceptions to the commutator-subgroup statement (for instance
$\GL_2(\F_2)\cong S_3$), the one-dimensional representations of $\GL_n(\F_q)$ are
$\chi\circ\det$ for $\chi$ a character of the cyclic group
$\F_q^\times\cong\Z_{q-1}$. The output determinant is a lazy random walk on $\Z_{q-1}$ with
steps $\{s_j=\mathrm{dlog}\det\sigma_j\}$.

\begin{proposition}[Determinant layer]
\label{prop:det}
\emph{(Strict contraction.)} The determinant layer has no non-trivial character of modulus one iff
$\{\det\sigma_j\}$ generate $\F_q^\times$. \emph{(Fourier coefficient.)} For the character
$\chi_t$ after $L$ leaves the complex coefficient is
$\prod_j(1+\chi_t(\det\sigma_j))/2$, and its modulus is
$\prod_j\abs{\cos(\pi t s_j/(q-1))}$. Here $L$ counts expanded occurrences, including inverse copies. No rate follows
from generation alone without specifying the multiset of steps and the tolerance.
For the separate lazy walk repeatedly using one generator of a cyclic group of
order $n\ge3$, mixing time at fixed tolerance in $(0,1/2)$ is $\Theta(n^2)$. \emph{(Leaves.)} An involution
$\sigma$ has $\det\sigma\in\{\pm1\}$; thus for $q-1>2$ involution leaves trap $\det$ in a
proper subgroup, and full $\GL_n(\F_q)$ mixing \emph{requires non-involution leaves}
whose determinants generate $\F_q^\times$. {\normalfont(In $\PGL_n$, the determinant class lies in
$\F_q^\times/(\F_q^\times)^n$, of order $\gcd(n,q-1)$; this is the square-class quotient only in
the essentially $n=2$ case.)}
\end{proposition}

The displayed Fourier coefficients quantify the determinant contribution
to mixing for the prescribed leaf multiset.

\paragraph{(c) The composite-modulus layer.}
For $q=p^k$, $1\to K\to\GL_n(\Z/p^k)\to\GL_n(\F_p)\to1$ with
$K=1+pM_n(\Z/p^k\Z)$ a nilpotent
$p$-group. Valid trees at arbitrary depth require a non-solvable ambient
group, and $\GL_n(\Z/p^k)$ is non-solvable iff
$\GL_n(\F_p)$ is. Since $\GL_2(\F_2)\cong S_3$ and $\GL_2(\F_3)$ are solvable, $\GL_2(\Z_4)$
and $\GL_2(\Z_9)$ are solvable: indeed $\GL_2(\Z_4)$ has derived series
$96\to24\to4\to1$, so every level-$3$ target is trivial.
For $k=2$, $K\cong M_n(\F_p)^{+}$ is abelian. Clifford theory then organises the irreducibles
above non-trivial $K$-characters by stabilisers, orbits, and induction.
The first non-degenerate $2\times2$ case is $\GL_2(\Z_{25})$, of order
$3\cdot10^5$. For generating conjugacy classes, Serre's lifting lemma
gives a direct reduction.

\begin{proposition}[Congruence lifting with generating determinants]
\label{prop:congruence-class}
Let $p\ge5$ be prime, $r\ge2$, $G=\GL_2(\Z/p^r)$, and
$K=\ker(G\to\GL_2(\F_p))$. Let $\mathcal C$ be conjugation-closed,
generate $G$, and have the property that $\det c$ generates
$(\Z/p^r)^\times$ for every $c\in\mathcal C$.
Then $(G,\mathcal C)$ has a uniform contraction level if and only if
$(\GL_2(\F_p),\bar{\mathcal C})$ does. In this case their exact levels
satisfy
\[
 \bar D\le D\le\bar D+\operatorname{dl}(K).
\]
In particular, $\bar D\le D\le\bar D+1$ for $r=2$.
The determinant hypothesis holds for every single conjugacy class
that generates $G$.
\end{proposition}
\begin{proof}
Every non-trivial $\mathcal C$-generated subgroup $H$ has full
determinant image. If its reduction is $\GL_2(\F_p)$, Serre's lifting
lemma~\cite[Lemma 3.15]{Lombardo2015}, applied to the closed preimage of
$H$ in $\GL_2(\Z_p)$, implies that $H$ contains $\SL_2(\Z/p^r)$.
The full determinant image then gives $H=G$. Consequently every proper
$\mathcal C$-generated subgroup has proper reduction. Since $K$ is
solvable, a non-solvable obstruction above yields one below; conversely,
\Cref{prop:quotient} lifts every obstruction below.
In the unobstructed case,
\[
 \operatorname{dl}(H)\le
 \operatorname{dl}(\bar H)+\operatorname{dl}(H\cap K)
 \le\bar D+\operatorname{dl}(K)
\]
for every proper $\mathcal C$-generated $H$. The opposite inequality
$\bar D\le D$ follows from \Cref{prop:quotient}.
For $r=2$, $K\cong M_2(\F_p)^+$ is abelian. Finally, all elements of
a single conjugacy class have the same determinant, which must generate
the unit group if the class generates $G$.
\end{proof}

\begin{remark}[A congruence obstruction for a union of classes]
\label{rem:congruence-supplement}
In $G=\GL_2(\Z/25)$ put $c=\operatorname{diag}(7,1)$ and
$\mathcal C=c^G\cup\{6I\}$. Then
$H=\langle c^G\rangle=\det^{-1}(\langle7\rangle)$ is a proper
non-solvable $\mathcal C$-generated subgroup of index $5$.
Indeed, conjugating $c$ by elementary unipotents and multiplying by
$c^{-1}$ gives all upper and lower unipotents: the respective
coefficients $1-7=-6$ and $1-7^{-1}=8$ are units modulo $25$.
They generate $\SL_2(\Z/25)$, proving the assertion about $H$.
The determinants $7$ and $\det(6I)=11$ have orders $4$ and $5$, so
$\mathcal C$ generates $G$.
Its reduction $\bar{\mathcal C}=\operatorname{diag}(2,1)^{\GL_2(5)}
\cup\{I\}$ is unobstructed. Every non-trivial
$\bar{\mathcal C}$-generated subgroup has full determinant image;
every non-solvable subgroup of $\GL_2(5)$ contains $\SL_2(5)$ by
Dickson's classification~\cite[Theorem 3.13]{Lombardo2015}, so these
two properties force it to be the whole group.
\end{remark}

\Cref{sec:rankone} gives the rank-one involution classification, followed
by the family rates and octahedral estimates in
\Cref{sec:quantitative-families,sec:octahedral-pair}.

%% ===================================================================
\section{The synthesis: two solvability conditions}
\label{sec:synthesis}
%% ===================================================================

By \Cref{thm:criterion}, strict contraction for a configuration is
equivalent to its leaves generating $G$, or escaping every maximal
subgroup. This places the matrix and character obstructions in a common
framework: each detects a proper leaf subgroup. For a conjugation-closed
generating leaf set in the configurable model, \Cref{thm:dichotomy}
identifies the non-degenerate mixing regime through two conditions:

\begin{enumerate}[nosep,label=(\roman*)]
\item the ambient group is \textbf{non-solvable}, so the commutator tree has a non-trivial
  accept value (a solvable ambient degenerates, as $\GL_2(\Z_4)$ shows); and
\item every \textbf{proper $\mathcal{C}$-generated subgroup is solvable}, with the level
  controlled by the largest derived length among them.
\end{enumerate}

These are exactly the two facts the companion paper~\cite{ZermenoS5} noted about $S_5$ in isolation
($S_5$ non-solvable; the point-stabiliser $S_4$ solvable of derived length $3$), here a
theorem for conjugation-closed leaf sets. \Cref{tab:synth} summarises the verified
evidence and the structural obstructions, and \Cref{cor:Sn,cor:A5,cor:pgl2} record the
exact statements.

\begin{table}[t]
\centering\small
\begin{tabular}{p{.27\textwidth}p{.40\textwidth}p{.22\textwidth}}
\toprule
Group and leaves & Decisive proper subgroups & Exact level / obstruction \\
\midrule
$S_5$, transpositions & $S_4$ has derived length $3$ & $3$ \\
$S_n$, transpositions, $n\ge6$ & $S_{n-1}$ non-solvable & Obstructed at every depth \\
$A_5$, involutions & $S_3,D_{10}$ attain derived length $2$ & $2$ \\
$\PSL_2(7)$, involutions & $S_4$; odd-order $7{:}3$ excluded & $3$ \\
$\PGL_2(q)$, all involutions, odd $q\ge5$ & $\PSL_2(q)$ & Obstructed \\
$\PGL_2(q)$, exterior involutions & Odd-degree non-solvable subfields; otherwise affine, dihedral or $S_4$ & $2$, $3$ or obstructed (\Cref{thm:rankone-classification}) \\
$\PSL_3(3)$, involutions & Parabolic $3^2{:}\GL_2(3)$ & $5$ \\
$\SL_3(3)$, transvections & $3^2{:}\SL_2(3)$ & $4$ \\
$\GL_2(\Z_4),\GL_2(\Z_9)$ & Ambient groups solvable & Degenerate \\
\bottomrule
\end{tabular}
\caption{Contraction levels and obstructions for the specified leaf sets.}
\label{tab:synth}
\end{table}

\begin{corollary}[Symmetric groups: obstruction at every depth]
\label{cor:Sn}
For $G=S_n$, $n\ge5$, with $\mathcal{C}$ the class of transpositions,
a uniform strict-contraction level exists exactly when $n=5$, and its exact
value is $3$ (in the configurable model; \Cref{rem:configurable}).
For every $n\ge6$ and every depth $m\ge0$, there is a valid configuration $\tau$
with all leaves in the point stabilizer $S_{n-1}<S_n$ whose operator satisfies
\[
 \opn{M_\tau(\rho_{\mathrm{std}})}=1,
\]
where $\rho_{\mathrm{std}}$ is the standard representation of $S_n$.
Consequently no depth gives strict contraction uniformly over all valid
configurations for any $n\ge6$.
\end{corollary}
\begin{proof}
A subgroup generated by transpositions is a Young subgroup $S_{a_1}\times\cdots\times
S_{a_r}$, corresponding to the connected components of the transposition
graph. For $n=5$, every proper such subgroup is solvable of derived length
at most $3$, and the point stabilizer $S_4$ attains $3$; hence
\Cref{thm:dichotomy} gives the exact level.

For $n\ge6$, let $H=S_{n-1}$ fix the last point. This is a proper
non-solvable subgroup generated by transpositions. By \Cref{lem:central},
the targets of configurations with leaves in the transposition class of
$H$ generate $H^{(m)}\ne1$ at every depth $m$. Thus at least one such
configuration is valid at each depth. Every leaf fixes
\[
 v=(1,\ldots,1,-(n-1))\ne0,
 \qquad v\in\{x\in\C^n:\textstyle\sum_i x_i=0\},
\]
so every leaf projection, and therefore the entire operator, fixes $v$.
As the operator is a product of contractions, its standard-representation
norm is exactly $1$. This gives the claimed obstruction at every depth.
\end{proof}

Exact shuffles by freely chosen lazy transpositions exist for every $S_n$
\cite{AngelHolroyd2018}, with suitable probabilities.
\Cref{cor:Sn} identifies when fair independent reads guarantee strict
contraction over all valid commutator trees.

\begin{corollary}[The alternating group $A_5$]
\label{cor:A5}
For $G=A_5$ with $\mathcal{C}$ the class of involutions, the construction strictly
contracts at exact level $2$ (configurable class; \Cref{rem:configurable}).
\end{corollary}
\begin{proof}
The proper $\mathcal{C}$-generated subgroups are $C_2$, $V_4$, $S_3$ and $D_{10}$---all
solvable, with $\max\operatorname{dl}=\operatorname{dl}(S_3)=\operatorname{dl}(D_{10})=2$.
(Note $A_4$ is \emph{not} $\mathcal{C}$-generated: its involutions generate only $V_4$,
illustrating \Cref{def:cgen}.) Hence $D=2$ by \Cref{thm:dichotomy}(2); an exhaustive check
confirms no level-$1$ configuration contracts, while all $29640$ ordered valid level-$2$ configurations
($3705$ reduced representatives) generate $A_5$. The numerical bottleneck is approximately $0.346$.
\end{proof}

\begin{corollary}[Spectral characterisation of minimal simple groups]
\label{cor:minsimple}
Let $G$ be a finite non-abelian simple group. The following are equivalent:
\begin{enumerate}[nosep,leftmargin=1.6em]
\item $G$ is minimal simple: every proper subgroup is solvable.
\item Every non-trivial conjugacy class admits a uniform strict-contraction
level in the configurable independent-read model.
\item Every non-trivial conjugacy class has uniform strict contraction at
level $5$.
\end{enumerate}
When these conditions hold, the exact level from \Cref{thm:dichotomy}
satisfies $D(G,\mathcal C)\le5$ for every non-trivial conjugacy class
$\mathcal C$.
The universal bound $5$ is attained by the involution class of
$\PSL_3(3)$ (\Cref{prop:psl33-level}).
\end{corollary}
\begin{proof}
Every non-trivial conjugacy class $\mathcal C$ generates $G$, since
$\langle\mathcal C\rangle$ is non-trivial and normal. If $G$ is minimal
simple, all proper $\mathcal C$-generated subgroups are solvable and have
derived length at most $5$~\cite[Remark 2.3]{ErsoyTortoraTota2014}.
\Cref{thm:dichotomy} gives the exact-level formula and (1)$\Rightarrow$(3).
The implication (3)$\Rightarrow$(2) is immediate.

For (2)$\Rightarrow$(1), suppose that $H<G$ is non-solvable. Choose
$x\in H\setminus O_\infty(H)$, where $O_\infty(H)$ is the largest soluble
normal subgroup of $H$, and set $\mathcal C=x^G$ and
$L=\langle\mathcal C\cap H\rangle$. Conjugation by $H$ preserves
$\mathcal C\cap H$, so $L\trianglelefteq H$. Since $x\in L$, solvability
of $L$ would contradict $x\notin O_\infty(H)$. Moreover $L\le H<G$ and
$\mathcal C\cap L=\mathcal C\cap H$, so $L$ is a proper non-solvable
$\mathcal C$-generated subgroup. \Cref{thm:dichotomy}(1) gives an
obstruction at every depth, contradicting (2).
\end{proof}

Thus contraction for every non-trivial class singles out precisely the
groups in Thompson's classification~\cite{Thompson1968}:
$\PSL_2(2^p)$ ($p$ prime), $\PSL_2(3^p)$ ($p$ odd prime), $\PSL_2(p)$ for
primes $p>3$ with $p\equiv\pm2\pmod5$, the Suzuki groups
$\mathrm{Sz}(2^p)$ ($p$ odd prime), and $\PSL_3(3)$.

\begin{remark}[$\GL_3(\F_2)$ has $D=3$ algebraically]
\label{rem:gl32}
$\GL_3(\F_2)\cong\PSL_2(7)$ is minimal simple, with maximal subgroups two classes of $S_4$
and one of $7{:}3$, all solvable. For the involution class, the proper
$\mathcal{C}$-generated subgroups have maximal derived length $\operatorname{dl}(S_4)=3$
(the $7{:}3$ is not involution-generated), so $D=3$ by \Cref{thm:dichotomy}(2). This is an
exact algebraic value; the level-$3$ enumeration in \Cref{sec:computation} only confirms it.
\end{remark}

\begin{corollary}[Projective linear groups]
\label{cor:pgl2}
Let $q$ be odd, $q\ge5$, and $\mathcal{C}$ the set of all involutions of $\PGL_2(\F_q)$
(a union of conjugacy classes, hence conjugation-closed). The involutions of
$\PSL_2(\F_q)$ form a single non-trivial conjugacy class of the simple group
$\PSL_2(\F_q)$ ($q\ge5$), hence generate it; adjoining any external involution---for
instance the image of $\left(\begin{smallmatrix}0&1\\ a&0\end{smallmatrix}\right)$ with
$-a$ a non-square, which squares to $aI$ (hence is an involution in $\PGL_2$) and has
determinant $-a$ (a non-square, hence external)---gives
$\langle\mathcal{C}\rangle=\PGL_2(\F_q)$. Thus the inner involutions generate the
proper non-solvable $\mathcal{C}$-generated subgroup $\PSL_2(\F_q)$ (index $2$), so by
\Cref{thm:dichotomy}(1) this leaf set has no uniform strict-contraction level. {\normalfont(Restricting $\mathcal{C}$ to the
$\PSL_2$-external involutions removes $\PSL_2$ from the $\mathcal{C}$-generated subgroups,
the exact positive and obstructed cases appear in \Cref{thm:rankone-classification}.)}
\end{corollary}

\subsection{Rigid templates: realisable levels and target persistence}
\label{subsec:rigid}
A rigid template fixes one leaf at each position. The per-configuration
criterion and the inclusion of level-$k$ targets in $H^{(k)}$ give the
upper bound $D$ for every rigid family (\Cref{prop:rigidupper}). We
characterise persistence of valid binary continuations through the
solvable radical, and show that persistence can be achieved by
conjugated doubling
(\Cref{thm:diagred,cor:binaryred,cor:binarylevel}). Over the transposition
class of $S_5$, the attainable non-vacuous exact levels are $2$ and $3$
(\Cref{thm:s5levels}).

\begin{definition}[Rigid families; non-vacuous exact level]
\label{def:rigidfamily}
A \emph{rigid family} over a leaf set $\mathcal{C}$ is a sequence
$F=(F_k)_{k\ge0}$ of specified sets of level-$k$ templates with leaves
in $\mathcal{C}$. Its \emph{formal threshold} $L(F)$ is the least integer
$d\ge0$ such that every valid $\tau\in F_k$, for every $k\ge d$, has
$H(\tau)=G$; set $L(F)=\infty$ if no such integer exists.
$F$ has \emph{non-vacuous exact level} $D_0\in\mathbb{Z}_{\ge0}$ if (i) $L(F)=D_0$; (ii) some
$\tau\in F_{D_0-1}$ is valid ($\alpha(\tau)\neq e$) with $H(\tau)<G$ (for $D_0=0$ this condition
is dropped); and (iii) every level $k\ge D_0$ contains a valid member. Condition (iii) rules out
vacuous thresholds: a family that is empty above one level may satisfy (i)--(ii) with no
content. By \Cref{thm:criterion}, $L(F)\le d$ says exactly that every valid member of level
$\ge d$ contracts strictly in every non-trivial irreducible.
\end{definition}

\begin{proposition}[Inherited upper bound]
\label{prop:rigidupper}
Let $\mathcal{C}$ be conjugation-closed, suppose every proper subgroup generated by a subset of
$\mathcal{C}$ is solvable, and let $D$ be the maximum derived length among those subgroups (the
mixing case of \Cref{thm:dichotomy}). Then $L(F)\le D$ for every rigid family $F$ over
$\mathcal{C}$.
\end{proposition}
\begin{proof}
Let $\tau\in F_k$ with $k\ge D$ and $\alpha(\tau)\neq e$. By \Cref{lem:valid},
$\alpha(\tau)\in H(\tau)^{(k)}$. If $H(\tau)<G$ then $H(\tau)$ is a proper subgroup generated by
a subset of $\mathcal{C}$, hence solvable of derived length at most $D\le k$, so
$H(\tau)^{(k)}=1$ and $\alpha(\tau)=e$, a contradiction. Hence $H(\tau)=G$.
\end{proof}

\begin{definition}[Target transition graph]
\label{def:targetgraph}
For a finite group $G$, let $\mathcal{T}(G)$ be the digraph on the non-trivial conjugacy classes
of $G$ with an edge $w^G\to(w')^G$ whenever $w'=[u,v]$ for some $u,v\in w^G$. Since
$[gwg^{-1},hwh^{-1}]$ is conjugate to $[w,(g^{-1}h)w(g^{-1}h)^{-1}]$, the graph is computed by
sweeping $[w,hwh^{-1}]$ over $h\in G$, one sweep per class.
\end{definition}

\begin{lemma}[Transfer]
\label{lem:classtransfer}
If $w^G\to(w')^G$ then for every $u\in w^G$ and $u'\in(w')^G$ there are explicit $g,h\in G$ with
$[gug^{-1},huh^{-1}]=u'$. In particular the relation is well defined on classes, and a self-loop
at $w^G$ is equivalent to the existence of a \emph{fixed} pair $(g,h)$ with
$[gwg^{-1},hwh^{-1}]=w$ literally.
\end{lemma}
\begin{proof}
Write $w'=[g_0wg_0^{-1},h_0wh_0^{-1}]$. If $u=sws^{-1}$, take
$g=g_0s^{-1}$ and $h=h_0s^{-1}$; then $[gug^{-1},huh^{-1}]=w'$.
Simultaneous conjugation carries $w'$ to any $u'\in(w')^G$.
For the loop statement: if $[w,h'wh'^{-1}]=twt^{-1}$ for some
$t$, conjugating by $t^{-1}$ gives $[t^{-1}wt,\;t^{-1}h'wh'^{-1}t]=w$, and both arguments are
conjugates of $w$.
\end{proof}

For a template $\tau$, write $g\cdot\tau$ for the template obtained by
conjugating each leaf label $s$ to $gsg^{-1}$. Its target is
$\alpha(g\cdot\tau)=g\alpha(\tau)g^{-1}$.

\begin{theorem}[Persistence criterion for conjugated doubling]
\label{thm:persistence}
Let $\mathcal{C}$ be conjugation-closed and let $\tau$ be a valid template with $H(\tau)=G$ and
target $w_0=\alpha(\tau)\neq e$. Consider continuations by \emph{conjugated doubling},
$\tau_{k+1}=[g_k\cdot\tau_k,h_k\cdot\tau_k]$.
The leaves stay in $\mathcal{C}$, and $H(\tau_{k+1})=G$ whenever $H(\tau_k)=G$.
A continuation with all members valid exists if and only if $w_0^G$ reaches a cycle of
$\mathcal{T}(G)$; the conjugator schedule may then be taken eventually periodic, with one fixed
pair per edge of the cycle.
\end{theorem}
\begin{proof}
Targets obey $\alpha(\tau_{k+1})=[\,g_k\alpha(\tau_k)g_k^{-1},\,h_k\alpha(\tau_k)h_k^{-1}\,]$, so
the target classes of a valid continuation trace an infinite walk in the finite digraph
$\mathcal{T}(G)$, which must reach a cycle. Conversely, fix a canonical representative of each
class on a path from $w_0^G$ into a cycle and around it; by \Cref{lem:classtransfer} each edge is
realised by an explicit conjugator pair carrying one canonical representative exactly to the
next, so both targets and conjugators are periodic from the moment the cycle is entered.
\end{proof}

We now allow binary compositions with children from different target classes. Fix a
conjugation-closed leaf set $\mathcal{C}$; all templates below are over $\mathcal{C}$, so
conjugated copies keep their leaves in $\mathcal{C}$. Call a rigid family
$(F_k)_{k\ge d}$ a \emph{binary continuation} of $F_d$ if every $\tau\in F_{k+1}$ has the form
$[g\cdot\tau_1,h\cdot\tau_2]$ with $\tau_1,\tau_2\in F_k$ (possibly equal) and $g,h\in
G$; call it \emph{persistently valid} if
$F_k\neq\emptyset$ \emph{and} every member is valid, for every $k\ge d$. Conjugated doubling is
the singleton case. For a set $S$ of non-trivial conjugacy classes let $\Phi(S)$ be the set of
non-trivial classes of commutators $[x,y]$ with $x,y\in\bigcup_{C\in S}C$; the diagonal slice
($x,y$ in one class) recovers exactly the edges of $\mathcal{T}(G)$.

\begin{lemma}[Mixed transfer]
\label{lem:mixedtransfer}
If $[x,y]=z\neq e$ with $x\in A$, $y\in B$, then for every $u\in A$, $v\in B$ and $z'\in z^G$
there are explicit $g,h\in G$ with $[\,gug^{-1},\,hvh^{-1}\,]=z'$: take $g=ca$, $h=cb$ with
$aua^{-1}=x$, $bvb^{-1}=y$, $czc^{-1}=z'$.
\end{lemma}

\begin{theorem}[Diagonal reduction]
\label{thm:diagred}
Let $S$ be a non-empty set of non-trivial conjugacy classes of $G$ and
$N(S)=\langle\bigcup_{C\in S}C\rangle$. The following are equivalent:
\textup{(1)} $\Phi^n(S)\neq\emptyset$ for all $n$;
\textup{(2)} $N(S)$ is non-solvable;
\textup{(3)} some $C\in S$ has non-solvable normal closure $\langle C\rangle$, equivalently
$C\not\subseteq O_\infty(G)$, the solvable radical;
\textup{(4)} some $C\in S$ reaches a cycle of $\mathcal{T}(G)$.
In particular a class reaches a cycle of $\mathcal{T}(G)$ if and only if its normal closure is
non-solvable.
\end{theorem}
\begin{proof}
\emph{(1)$\Leftrightarrow$(2).} $\mathcal{C}_0=\bigcup_{C\in S}C$ is conjugation-closed, so
\Cref{lem:central} gives $\langle T_k\rangle=N(S)^{(k)}$ for $T_0=\mathcal{C}_0$,
$T_{k+1}=\{[a,b]:a,b\in T_k\}$. A commutator with a trivial argument is trivial, and each $T_k$
is conjugation-closed, so induction gives $T_k\setminus\{e\}=\bigcup_{C\in\Phi^k(S)}C$; hence
$\Phi^k(S)=\emptyset\iff N(S)^{(k)}=\{e\}$, and (1) holds iff $N(S)$ is non-solvable.
\emph{(2)$\Leftrightarrow$(3).} Each $\langle C\rangle$ is normal (generated by a normal
subset) and $N(S)=\prod_{C\in S}\langle C\rangle$; a product of finitely many normal solvable
subgroups is solvable, so non-solvability passes to some factor, and conversely each factor
lies in $N(S)$. Normality gives $\langle C\rangle$ solvable $\iff C\subseteq O_\infty(G)$.
\emph{(3)$\Rightarrow$(4).} If $N_0=\langle C_0\rangle$ is non-solvable then so is $N_0'$, and
\Cref{lem:central} (one level, with the class $C_0$) writes $N_0'=\prod\langle D\rangle$ over
the successors $C_0\to D$ in $\mathcal{T}(G)$; some successor has non-solvable closure, and
iterating drives an infinite walk of the finite digraph into a cycle.
\emph{(4)$\Rightarrow$(1).} An infinite walk $C_0\to C_1\to\cdots$ from $C_0\in S$ gives
$C_n\in\Phi^n(\{C_0\})\subseteq\Phi^n(S)$ by monotonicity.
\end{proof}

\begin{corollary}[Binary continuations reduce to doubling]
\label{cor:binaryred}
Let $F_d\neq\emptyset$ be a set of valid level-$d$ templates over $\mathcal{C}$ with
target-class set $S_0$. The
following are equivalent: \textup{(a)} some binary continuation of $F_d$ is persistently valid;
\textup{(b)} $\Phi^n(S_0)\neq\emptyset$ for all $n$; \textup{(c)} some $\tau\in F_d$ has target
class $\not\subseteq O_\infty(G)$; \textup{(d)} some single $\tau\in F_d$ admits a valid
conjugated-doubling continuation. Mixed-class children never extend the reach of persistence.
\end{corollary}
\begin{proof}
(a)$\Rightarrow$(b): the target of $[g\cdot\tau_1,h\cdot\tau_2]$ is
$[\,g\alpha(\tau_1)g^{-1},h\alpha(\tau_2)h^{-1}\,]$, so target classes of valid level-$(d+n)$
members lie in $\Phi^n(S_0)$ by induction. (b)$\Rightarrow$(a): maintain one valid member per
class of $\Phi^n(S_0)$, composing by \Cref{lem:mixedtransfer} (the same member, doubly
conjugated, when the two classes coincide): the children are conjugated copies of valid
templates and the root target is a prescribed $z'\neq e$. (b)$\Leftrightarrow$(c):
\Cref{thm:diagred}. (c)$\Rightarrow$(d): by \Cref{thm:diagred} the target class reaches a
cycle, and the realisation in the proof of \Cref{thm:persistence} uses only the target walk.
(d)$\Rightarrow$(a): doubling is a binary continuation.
\end{proof}

\begin{corollary}[Prefixes admitting an exact-$D$ binary continuation]
\label{cor:binarylevel}
In the mixing case of \Cref{thm:dichotomy} with $D\ge1$: given $F_0,\dots,F_{D-1}$ and a
non-empty set $F_D$ of valid level-$D$ templates, some binary continuation of $F_D$ yields a
rigid family of non-vacuous exact level $D$ if and only if \textup{($\alpha$)} some
$\tau^*\in F_{D-1}$ is valid with $H(\tau^*)<G$, and \textup{($\beta$)} some $\tau\in F_D$ has
target class $\not\subseteq O_\infty(G)$. The continuation may then be taken to be conjugated
doubling of a single member of $F_D$; when $O_\infty(G)=\{e\}$, condition \textup{($\beta$)}
holds for every valid member. This characterises binary \emph{extendibility} of a prefix;
classifying \emph{prescribed} families remains open (\Cref{sec:open}).
\end{corollary}
\begin{proof}
$L(F)\le D$ by \Cref{prop:rigidupper}, whose argument shows every valid member at level $\ge D$
has $H=G$; hence conditions (i)--(ii) of \Cref{def:rigidfamily} reduce to ($\alpha$) (and
$\operatorname{dl}(H(\tau^*))=D$ automatically, by validity at level $D-1$). Condition (iii)
holds for some binary continuation iff ($\beta$): sufficiency is \Cref{cor:binaryred}; for
necessity pass to the subfamily of valid members---subtrees of a valid template are valid, so
it is again a binary continuation, persistently valid by (iii)---and apply
\Cref{cor:binaryred}. The graph data, the transfer identities, the subset-exhaustive
equivalence of \Cref{thm:diagred} over nine groups, and a five-level $S_5$ family with
mixed-class children at every production are certified in the accompanying suite.
\end{proof}

\begin{theorem}[Realisable levels over $S_5$]
\label{thm:s5levels}
Over the transposition class of $S_5$, the non-vacuous exact levels realisable by rigid families
are exactly $\{2,3\}$.
\end{theorem}
\begin{proof}
Exact level $0$ is impossible: persistence would require a valid member of $F_0$, i.e.\ a single
leaf, which generates a group of order $2<|S_5|$, so $L(F)\ge1$. Exact level $1$ is impossible:
persistence would require a valid member of $F_1$, and a valid pair of transpositions generates
$S_3<S_5$, so $L(F)\ge2$. And $L(F)\le3$ always, by \Cref{prop:rigidupper} with $D=3$
(\Cref{thm:dichotomy}, \Cref{cor:Sn}).

Level $2$ is realised by the family $F$ with $F_0=\{(12),(23),(34),(45)\}$,
$F_1=\{[(12),(23)],\allowbreak[(34),(45)]\}$ (the leaf configurations and level-$1$ subtrees of
the seed), and $F_k=\{\gamma_k\}$ for $k\ge2$, where $\gamma_2=[[(12),(23)],[(34),(45)]]$ and
$\gamma_{k+1}=[g\cdot\gamma_k,h\cdot\gamma_k]$ with the fixed pair $g=(4\,5)$,
$h=(1\,2)(3\,5)$. The seed generates $S_5$ with target the $3$-cycle $w_3=(2\,5\,3)$, and
$[gw_3g^{-1},hw_3h^{-1}]=w_3$ (a self-loop of $\mathcal{T}(S_5)$ at the $3$-cycle
class), so by induction every $\gamma_k$, $k\ge2$, is valid with target $w_3$ and generates
$S_5$. Hence $L(F)=2$ (the defining condition holds at $d=2$ and fails at $d=1$ through the
$S_3$ members of $F_1$), the level-$1$ witness is either member of $F_1$, and valid members
exist at every level $\ge2$: non-vacuous exact level $2$.

Level $3$ is realised by the family $E$ with $E_0,E_1$ the leaf configurations and level-$1$
subtrees of the seed $\tau_3$ below, $E_2=\{\tau_A,\tau_B\}$ for
$\tau_A=[[(12),(23)],[(34),(14)]]$ and $\tau_B=[[(23),(35)],[(45),(24)]]$, and
$E_k=\{\tau_k\}$ for $k\ge3$, where $\tau_3=[\tau_A,\tau_B]$ and
$\tau_{k+1}=[g'\cdot\tau_k,h'\cdot\tau_k]$ with the fixed pair $g'=(3\,4\,5)$,
$h'=(2\,3\,4)$. Each of $\tau_A,\tau_B$ is valid and generates a point stabiliser $S_4<S_5$;
$\tau_3$ generates $S_5$ with target the $5$-cycle $w_5=(1\,3\,2\,4\,5)$, and
$[g'w_5g'^{-1},h'w_5h'^{-1}]=w_5$, so by induction every $\tau_k$, $k\ge3$, is valid with
target $w_5$ and generates $S_5$. Hence $L(E)=3$ (the condition holds at $d=3$ and fails at
$d=2$ through $\tau_A$), the level-$2$ witness is $\tau_A$, and valid members exist at every
level $\ge3$: non-vacuous exact level $3$.

The accompanying suite certifies these identities and generation facts,
and verifies the resulting families through level $6$.
\end{proof}

\begin{remark}[Connection with the companion paper]
\label{rem:p0scope}
The persistence criterion supplies valid continuations, and
\Cref{thm:criterion} gives qualitative contraction when their leaves
generate $G$. The companion paper~\cite{ZermenoS5} provides a quantitative
mixing estimate for a specified fixed-template point-function construction.
\Cref{thm:s5levels} classifies the attainable exact levels across rigid
families over the transposition class.
\end{remark}

%% ===================================================================
\section{Involution classes in rank one}
\label{sec:rankone}

Here the dichotomy gives an explicit classification for an infinite family.
Throughout this section $q=p^f\ge5$ is odd. In $\PGL_2(q)$ the determinant modulo
squares defines a character with kernel $\PSL_2(q)$. Call an involution
\emph{interior} or \emph{exterior} according as its determinant is a square or a
non-square. These are the two involution classes of $\PGL_2(q)$.
The subgroup classification used below is Dickson's, in the form recorded by
Giudici~\cite[Theorems 2.2 and 3.5]{Giudici2007}; our conclusions concern the
commutator-tree levels for these specified leaf sets.

\begin{lemma}[Classes and subfields]
\label{lem:rankone-inputs}
The interior involutions generate $\PSL_2(q)$, and the exterior involutions
generate $\PGL_2(q)$. Every $A_5$ subgroup of $\PGL_2(q)$ is interior.
For a proper subfield $\F_{q_0}\subset\F_q$, exterior involutions of
$\PGL_2(q_0)$ remain exterior precisely when $[\F_q:\F_{q_0}]$ is odd.
\end{lemma}
\begin{proof}
The involutions of the simple group $\PSL_2(q)$ form a non-trivial normal
subset, so generate it. The subgroup generated by the exterior class is
normal in $\PGL_2(q)$. Its intersection with the simple socle is either the
socle or trivial; in the latter case its commutator with the socle is trivial,
contradicting the trivial centralizer of the socle in $\PGL_2(q)$. Hence it
contains the socle and an exterior element, and is the whole group.
The determinant character is trivial on the perfect group $A_5$.
Finally, if $m=[\F_q:\F_{q_0}]$, the parity of
$(q-1)/(q_0-1)=1+q_0+\cdots+q_0^{m-1}$ is the parity of $m$.
The square-character restriction therefore has the stated parity.
\end{proof}

\begin{theorem}[Rank-one classification]
\label{thm:rankone-classification}
For odd $q\ge5$ the following hold.
\begin{enumerate}[nosep,leftmargin=1.6em]
\item In $G=\PGL_2(q)$, the set of \emph{all} involutions has a
permanent obstruction.
\item In $G=\PGL_2(q)$, the exterior class has a permanent obstruction if and
only if there is a proper subfield $\F_{q_0}$ with $q_0\ge5$ and
$[\F_q:\F_{q_0}]$ odd. Otherwise its exact level is
\[
D=\begin{cases}3,&q\equiv\pm3\pmod8,\\2,&q\equiv\pm1\pmod8.\end{cases}
\]
\item In $G=\PSL_2(q)$, the involution class has a permanent obstruction if
and only if $G$ contains a \emph{proper} subgroup isomorphic to $A_5$, or
$\F_q$ has a proper subfield of size $q_0\ge5$. Otherwise its exact level is
\[
D=\begin{cases}3,&q\equiv\pm1\pmod8,\\2,&q\equiv\pm3\pmod8.\end{cases}
\]
In particular $\PSL_2(5)=A_5$ has level $2$. The $A_5$ obstruction in this
part occurs exactly when $q>5$ and either $p=5$ or $q\equiv\pm1\pmod5$.
\end{enumerate}
\end{theorem}
\begin{proof}
For (1), the interior leaves generate the proper non-solvable subgroup
$\PSL_2(q)$, so \Cref{thm:dichotomy} applies.

For (2), a proper exterior-generated subgroup lies in a maximal subgroup
other than $\PSL_2(q)$. Giudici's Theorem~3.5 lists these maximals: an affine
group $\F_q^+\rtimes\F_q^\times$, torus normalizers (dihedral groups), $S_4$
when $q=p\equiv\pm3\pmod8$, and $\PGL_2(q_0)$ when $q=q_0^r$ with $r$ an odd
prime. The affine group is metabelian, including over extension fields;
its involution-generated subgroups need not be ordinary dihedral groups.
Torus normalizers are also metabelian. The only non-solvable possibility
outside the socle is therefore a proper subfield group with $q_0\ge5$ and
odd degree. Such a group is indeed exterior-generated by
\Cref{lem:rankone-inputs}. Any odd proper degree has an odd prime divisor,
so the existence criterion agrees with this maximal-subgroup list.

In its absence, the possible subfield maximal with $q_0=3$ is
$\PGL_2(3)=S_4$, occurring for $q=3^r$, $r$ an odd prime. Together with the
prime-field $S_4$ case, these give an exterior-generated $S_4$ exactly in
the unobstructed cases $q\equiv\pm3\pmod8$.
(Indeed, this congruence forces $f$ odd; without a proper subfield of size
at least $5$ and odd degree, $f=1$ when $p\ge5$, and $f$ is an odd prime
when $p=3$.) Its determinant character is
the sign character, so its transpositions are exterior and its derived
length $3$ is attained. In the remaining cases all relevant maximals are
metabelian. The upper bound $2$ is attained: for $q\ge7$, take a split
torus generator $r:x\mapsto ax$ and an exterior involution $t:x\mapsto b/x$
($-b$ non-square). The two exterior involutions $t,r^2t$ generate
$\langle r^2,t\rangle$, a proper non-abelian dihedral group whose rotation
subgroup has order $(q-1)/2\ge3$. The case $q=5$ has the exterior $S_4$ and
was already covered. Apply \Cref{thm:dichotomy} for exactness.

For (3), Dickson's list says that any proper non-solvable subgroup lies in
an $A_5$ or a non-solvable subfield subgroup. A proper $A_5$ is generated
by its involutions. Every proper subfield of size $q_0\ge5$ supplies a
proper involution-generated $\PSL_2(q_0)$, irrespective of the degree;
even degree can also supply $\PGL_2(q_0)$. Conversely, the absence of these
two obstructions leaves affine and dihedral subgroups of derived length
at most $2$, and $A_4$ or $S_4$. The involutions in $A_4$ generate only
$V_4$. An $S_4$ occurs precisely for $q\equiv\pm1\pmod8$ and is generated
by involutions, giving level $3$ when unobstructed. Otherwise a
non-abelian dihedral torus normalizer gives level $2$: at least one of
$(q-1)/2,(q+1)/2$ is at least $3$, and its dihedral normalizer is proper.
This argument uses existence, not maximality, of these dihedral groups
at small $q$. At $q=5$ the $A_5$ is the ambient group and is not an
obstruction, in agreement with \Cref{cor:A5}. The stated $A_5$ embedding
condition is the corresponding part of Dickson's classification.
\end{proof}

\begin{corollary}[Prime fields]
\label{cor:rankone-prime}
For every prime $p\ge5$, the exterior involution class of $\PGL_2(p)$ mixes
at exact level $3$ for $p\equiv\pm3\pmod8$ and at exact level $2$ for
$p\equiv\pm1\pmod8$.
\end{corollary}

Thus $q=5,7,11$ give exterior levels $3,2,3$. At $q=11$ the involutions
of the $\PSL_2(11)$ ambient are instead obstructed by a proper $A_5$.
The extension degree matters: $q=25$ has exterior level $2$, whereas
$q=125$ contains the exterior-generated non-solvable $\PGL_2(5)$ and is
obstructed. \Cref{sec:quantitative-families} gives mixing rates as $q$ grows.


\section{Uniform rates and the order of the mixing depth}
\label{sec:quantitative-families}

The dichotomy supplies generating sets at bounded depth. Diameter estimates
then control their contraction uniformly over all valid configurations.
In particular, the unobstructed exterior
classes of $\PGL_2(q)$ have base depth $2$ or $3$, uniformly in $q$;
the issue here is the strength of contraction as the field grows.
All logarithms in this section are natural
unless their base is displayed. For a tree $\tau$, the law $\mu_\tau$
is unconditional.

\begin{theorem}[Diameter-to-mixing transfer]
\label{thm:family-rate}
Consider a family of pairs $(G,\mathcal C)$ satisfying case~\textup{(2)} of
\Cref{thm:dichotomy}, with exact levels $D\le D_0$. Suppose there are
constants $K,a>0$, independent of the pair, such that every symmetric
generating set $S$ of $G$ satisfies
$\mathrm{diam}(G,S)\le K(\log\abs G)^a$.
Put $c=(8\cdot16^{D_0}K^2)^{-1}$. Every valid tree $\tau$ of
level $m\ge D$ satisfies
\begin{equation}\label{eq:family-rate}
 \dtv(\mu_\tau,U_G)\le\frac{\sqrt{\abs G-1}}2
       \exp\left(-\frac{c4^m}{(\log\abs G)^{2a}}\right).
\end{equation}
For $D_T$, add the conditioning probability $p$ from \Cref{lem:conditioning}.
The leaves need not be involutions.
\end{theorem}
\begin{proof}
Every sub-block of a valid tree is valid. Its level-$D$ sub-blocks therefore
generate $G$ by \Cref{thm:dichotomy}. Each has $N=4^D$ expanded occurrences,
including inverse occurrences. Fix a non-trivial irreducible $\rho$.
For each level-$D$ sub-block $\tau_0$, \Cref{prop:stability} gives
\[
 \opn{M_{\tau_0}(\rho)}\le1-\frac1{8\cdot4^D K^2(\log\abs G)^{2a}}.
\]
Expanding the recursion down to this level produces $4^{m-D}$ factors,
counting adjoints, all with that norm bound. Submultiplicativity and
$1-x\le e^{-x}$ yield
\[
 \opn{M_\tau(\rho)}\le
 \exp\left(-\frac{4^m}{8\cdot16^D K^2(\log\abs G)^{2a}}\right).
\]
Apply Diaconis--Shahshahani as in \Cref{prop:quant}, and use $D\le D_0$.
\end{proof}

For $0<\eta<1/2$, a sufficient condition for error at most $\eta$ is
$m\ge D$ and
\begin{equation}\label{eq:family-accuracy}
 4^m\ge c^{-1}(\log\abs G)^{2a}
       \left(\tfrac12\log(\abs G-1)+\log\frac1{2\eta}\right).
\end{equation}
Thus fixed error, or error $\abs G^{-b}$ for fixed $b>0$, requires only
$O(\log\log\abs G)$ depth and $O((\log\abs G)^{2a+1})$ expanded length.

\begin{proposition}[Support lower bound and optimal order]
\label{prop:depth-lower}
For any level-$m$ independent-read tree $\tau$,
\[
 \dtv(\mu_\tau,U_G)\ge\max\{0,1-2^{4^m}/\abs G\}.
\]
Consequently, for any fixed $0<\eta<1/2$, in an unbounded family satisfying
\Cref{thm:family-rate}, the least depth at which every valid tree has
error at most $\eta$ is $\Theta(\log\log\abs G)$, uniformly over the family.
\end{proposition}
\begin{proof}
There are $4^m$ independent input bits, hence at most $2^{4^m}$ possible
outputs. The complement of their support has uniform mass at least
$1-2^{4^m}/\abs G$ and output mass zero. Error at most $\eta$ therefore
requires $4^m\ge\log_2((1-\eta)\abs G)$.
Valid trees exist at every level by \Cref{lem:central}, since $G$ is
non-solvable. This rules out vacuous small levels. The upper bound is
\eqref{eq:family-accuracy}.
\end{proof}

The upper and lower bounds match in their order of depth, up to constant factors.
On the obstructed side, the subgroup $H<G$ in \Cref{thm:dichotomy}(1)
supplies at every level a valid tree $\tau$ supported in $H$, and hence
$\dtv(\mu_\tau,U_G)\ge1-\abs H/\abs G\ge1/2$.

\begin{corollary}[Rank one and minimal simple groups]
\label{cor:family-examples}
The uniform rate \eqref{eq:family-rate} and the optimal order of depth in
\Cref{prop:depth-lower} apply to:
\begin{enumerate}[nosep,leftmargin=1.6em]
\item the unobstructed involution classes of $\PSL_2(q)$ and exterior
classes of $\PGL_2(q)$ in \Cref{thm:rankone-classification};
\item all finite minimal simple groups with any non-trivial conjugacy
class, uniformly over both the group and the class.
\end{enumerate}
\end{corollary}
\begin{proof}
Pyber--Szab\'o's bounded-rank diameter theorem~\cite[Theorem 2 of the
preprint]{PyberSzabo2016} supplies constants $K,a$ for every symmetric
generating set of a simple group of Lie type of bounded rank.
For $G=\PGL_2(q)$, take a symmetric generating set $S$ and $t\in S$
outside $L=\PSL_2(q)$. With transversal $\{1,t\}$, Schreier generators
of $L$ have length at most $3$ in $S$; their inverses satisfy the same
bound. Writing this symmetric set as $S_L$ gives
\[
 \mathrm{diam}(G,S)\le3\,\mathrm{diam}(L,S_L)+1.
\]
Thus the same type of diameter bound holds for $G$, with adjusted absolute
constants. The rank-one levels are at most $3$.
For (2), Thompson's list in \Cref{cor:minsimple} has bounded Lie rank
(at most $2$, including the underlying type of the Suzuki groups), and
\Cref{cor:minsimple} gives $D\le5$. Apply \Cref{thm:family-rate} with
$D_0=3$ or $5$, respectively, and then \Cref{prop:depth-lower}.
\end{proof}

\subsection{Uniform expansion outside exceptional primes}

Over prime fields, strong uniform expansion gives constant gaps for every
generating block of bounded depth. The prime-field input goes back to
Breuillard--Gamburd~\cite[Corollary 1.1]{BreuillardGamburd2010}.
Becker--Breuillard~\cite[Theorem 1.1 and Corollary 3.2]{BeckerBreuillard2026}
extend uniform expansion to bounded-rank families with bounded field degree,
outside exceptional sets of characteristics.

\begin{corollary}[Uniform gaps for almost all primes]
\label{cor:almost-all-primes}
There is a set $\mathcal E$ of relative density zero among the primes
such that the following hold for every prime $p\ge5$ outside $\mathcal E$.
\begin{enumerate}[nosep,leftmargin=1.6em]
\item For each fixed $K\ge1$ there is $a_K>0$, independent of $p$, such that
every involutive block $\tau$ of level $k\le K$ whose leaves generate
$G=\PSL_2(p)$ or $\PGL_2(p)$ satisfies
$\opn{M_\tau(\rho)}\le1-a_K$ for every non-trivial irreducible $\rho$.
\item Let $\mathcal C$ be an unobstructed involution class from
\Cref{thm:rankone-classification}: the exterior class in $\PGL_2(p)$,
or the involution class in an unobstructed $\PSL_2(p)$.
There is an absolute $c>0$ such that every valid tree $\tau$ of level $m\ge D$
over $\mathcal C$ satisfies
\begin{equation}\label{eq:almost-all-prime-rate}
 \opn{M_\tau(\rho)}\le e^{-c4^m},\qquad
 \dtv(\mu_\tau,U_G)\le\frac{\sqrt{\abs G-1}}2 e^{-c4^m}.
\end{equation}
The operator bound holds for every non-trivial irreducible $\rho$.
\item If $p\equiv\pm3\pmod8$, every pair of exterior subgroups
$A,B\cong S_4$ generating $\PGL_2(p)$ satisfies
$\opn{P_A^\rho P_B^\rho}\le1-b$ for every non-trivial irreducible $\rho$,
with an absolute $b>0$ independent of the pair and of $p$.
\end{enumerate}
\end{corollary}
\begin{proof}
Strong uniform expansion supplies $\varepsilon\in(0,1]$ and
$\mathcal E$ such that $|SX|\ge(1+\varepsilon)|X|$ whenever
$|X|\le|G|/2$ and $S$ is a symmetric generating set containing $e$,
for both groups $G$. Starting with
\cite[Corollary 1.1]{BreuillardGamburd2010}, pass from $\SL_2(p)$ to
its quotient $\PSL_2(p)$ and then to its index-two overgroup
$\PGL_2(p)$; uniform expansion is preserved under these operations
\cite[Lemma 3.9]{BeckerBreuillard2026}. Both sources use vertex
expansion: the proof of \cite[Corollary 1.1]{BreuillardGamburd2010}
produces a single $s\in S$ with $|sX\setminus X|\ge c|X|$, and
\cite[(1.1)]{BeckerBreuillard2026} states the displayed form directly.

For (1), fix such a block $\tau$ and list its $r=2^k$ leaf positions with their multiplicities.
Their lazy walk assigns each distinct non-identity leaf probability at
least $1/(2r)$. The expansion inequality gives at least
$\varepsilon|X|$ outgoing edges from any $X$ with $|X|\le|G|/2$,
so its conductance is at least $\varepsilon/(2r)$.
Cheeger's inequality gives spectral gap
$\delta\ge\varepsilon^2/(8r^2)$, including these multiplicities.
Each leaf occurs $2^k$ times in the expanded involutive block, so this
is also its expanded-leaf walk. Put $N=4^k=r^2$. By
\Cref{prop:cayley-comparison},
\[
 1-\opn{M_\tau(\rho)}^2
 \ge\frac{\varepsilon^2}{4N(N+1)}.
\]
Thus one may take $a_K=\varepsilon^2/(8\cdot4^K(4^K+1))$.

For (2), $D\le3$, and every level-$D$ sub-block of a valid tree $\tau$
generates $G$ by \Cref{thm:dichotomy}. Applying (1) at this fixed
base level and expanding the recursion gives
\[
 \opn{M_\tau(\rho)}\le(1-a_3)^{4^{m-D}}
 \le e^{-(a_3/64)4^m}.
\]
Take $c=a_3/64$ and apply the Diaconis--Shahshahani bound.

For (3), fix a non-trivial irreducible $\rho$ on $V$, and write
$P_A=P_A^\rho$, $P_B=P_B^\rho$.
The walk $\nu=(U_A+U_B)/2$ has generating support
$A\cup B$ and assigns each support element probability at least $1/48$.
The same conductance argument gives spectral gap at least
$\eta=\varepsilon^2/(2\cdot48^2)$. Its Fourier operator is
$(P_A+P_B)/2$, which is positive semidefinite.
For a unit vector $v\in V^A$,
$1+\norm{P_Bv}^2\le2(1-\eta)$.
Hence $\opn{P_AP_B}\le\sqrt{1-2\eta}\le1-\eta$;
if $V^A=0$ the norm is zero. Take $b=\eta$.
\end{proof}

Part~(3) is a qualitative form of \Cref{conj:allpairs} outside
$\mathcal E$: a gap for all generating exterior pairs, without the sharp
constant of \Cref{thm:octahedral-sharp}.

\subsection{Few exceptional configurations over arbitrary fields}

\begin{corollary}[Exceptional generating blocks at fixed depth]
\label{cor:few-exceptional-blocks}
Fix integers $R,k\ge1$ and $\delta>0$.
There is $a=a(R,k,\delta)>0$ such that, for every finite simple group
$G(q)$ of Lie type of rank at most $R$, all but at most
$O_{R,k,\delta}(q^\delta)$ simultaneous $G(q)$-conjugacy classes of
generating involutive level-$k$ configurations $\tau$ satisfy
\[
 \opn{M_\tau(\rho)}\le1-a
 \qquad\text{for every non-trivial irreducible }\rho.
\]
The same conclusion holds for $\PGL_2(q)$, $q\ge4$.
Restricting to valid configurations, or to specified conjugacy classes
of involution leaves, preserves this absolute bound on exceptions.
\end{corollary}
\begin{proof}
The dimension-zero theorem of Becker--Breuillard
\cite[Theorems 1.2 and 4.1]{BeckerBreuillard2026}, with $2^k$ entries,
gives an expansion constant $\varepsilon=\varepsilon(R,k,\delta)>0$
outside the stated number of simultaneous conjugacy classes of
generating tuples. The finite simple groups have the required
bounded-dimensional algebraic embeddings; $\PGL_2(q)$ is covered
by Theorem~4.1 in the adjoint group of type $A_1$, since it contains
$\PSL_2(q)$ and therefore satisfies the density hypothesis of that
theorem for $q$ large.
Finitely many small fields are absorbed into the implied constant.
For an involutive tuple put $N=4^k$. Its lazy leaf walk, with
multiplicities, has gap at least $\varepsilon^2/(8N)$ by the
conductance argument in \Cref{cor:almost-all-primes}. By
\Cref{prop:cayley-comparison} one may take
$a=\varepsilon^2/(8N(N+1))$, after reducing $\varepsilon$ to at most $1$.
Conjugating every leaf conjugates the block operator, so the transfer
respects exactly the equivalence relation counted by the source theorem.
\end{proof}

Here the depth is fixed and the field degree is unrestricted. The
exception count is absolute; a relative-density statement in a
restricted family additionally requires a lower bound on its size.

\begin{remark}[A question]
\label{rem:rates-question}
Do the involution classes of the minimal simple families
$\PSL_2(2^\ell)$, $\PSL_2(3^\ell)$ and $\mathrm{Sz}(2^\ell)$ admit a
constant gap for every valid configuration of level $D$, uniformly as
the prime exponent $\ell$ grows? \Cref{cor:few-exceptional-blocks}
allows exceptional configurations; removing them in these families
would follow from the bounded-rank uniform expansion conjecture
discussed in~\cite[Introduction]{BeckerBreuillard2026}.
\end{remark}


\section{The natural octahedral pair: a sharp uniform angle}
\label{sec:octahedral-pair}

We construct a pair of octahedral subgroups whose fixed-space angle is
controlled by the Ramanujan bound for a quaternionic Cayley walk, and
show that it is the only pair of its kind up to conjugacy. A compression
identity connects the two operators.

\subsection{Fixed-space angles and double cosets}
\label{subsec:hecke-boundary}
Let $H$ be a subgroup of a finite group $G$ and let $\rho$ be a unitary
representation of $G$ on $V$. Put
\[
  e_H=\frac1{\abs H}\sum_{h\in H}h\in\C[G],
  \qquad
  P_H=P_H^\rho=\rho(e_H)=\frac1{\abs H}\sum_{h\in H}\rho(h).
\]
When $\rho$ is fixed we omit the superscript; we retain it when taking
maxima over representations.

\begin{lemma}[Double cosets and angles between fixed spaces]
\label{cor:double-coset-angle}
The operator $P_H$ is the orthogonal projection onto $V^H$. For $g\in G$
and $K=gHg^{-1}$,
\[
  P_K=\rho(g)P_H\rho(g)^{-1},
  \qquad
  \opn{P_H\rho(g)P_H}=\opn{P_HP_K}.
\]
The element $e_Hge_H$ depends only on the double coset $HgH$.
\end{lemma}
\begin{proof}
The average $P_H$ is self-adjoint and idempotent. Its image is $H$-fixed
by reindexing the sum, and it acts as the identity on every $H$-fixed
vector. Conjugating the average over $H$ gives the formula for $P_K$, so
$P_H\rho(g)P_H=P_HP_K\rho(g)$, and right multiplication by the unitary
$\rho(g)$ preserves the norm. Finally $e_Hh_1gh_2e_H=e_Hge_H$ for
$h_1,h_2\in H$.
\end{proof}

Thus $\opn{P_H\rho(g)P_H}$, the norm in $\rho$ of the element of the
Hecke algebra $e_H\C[G]e_H$ attached to the double coset $HgH$, is the
cosine of the smallest principal angle between $V^H$ and $V^K$, with
cosine $0$ when either space is zero. If $\langle H,K\rangle=G$ and
$\rho$ is non-trivial and irreducible, then $V^H\cap V^K=V^G=0$, and
this is the Friedrichs angle, which governs the alternation of the two
projections~\cite{KayalarWeinert1988,BadeaGrivauxMuller2010}.

\subsection{The pair and its four neighbours}
For the rest of this section $q$ is a prime with $q\equiv\pm3\pmod8$,
$q\ne3$, and $G_q=\PGL_2(\F_q)$.
Let $\mathbb H=(-1,-1)_{\mathbb Q}$, with $i^2=j^2=k^2=ijk=-1$,
and take projective classes in $\mathbb H^\times/\mathbb Q^\times$:
\[
 A_0=\langle i,1+i,1+i+j+k\rangle,\qquad h_0=i+j+k,
 \qquad S_0=\{\pm i\pm j\pm k\}/\{\pm1\}.
\]
The four elements of $S_0$ have norm $3$ and are projective involutions.
For any odd prime $q$, the quaternion algebra over $\F_q$ is split.
Indeed, the two sets of squares and of $-1$ minus squares have more
than $q$ elements in total, so $x^2+y^2=-1$ has a solution; this gives
an explicit splitting into $M_2(\F_q)$~\cite[Proposition 2.5.2]{DSV2003}.
Reducing the indicated projective elements defines $A,h,S$ in $G_q$.
Their norms $1,2,4,3$ are invertible for our primes. Changing the splitting
conjugates the entire construction, leaving the norms below unchanged.

\begin{lemma}[Octahedral reduction]
\label{lem:octahedral-construction}
One has $A\cong S_4$, with exterior transpositions, and
\[
 B=hAh^{-1},\quad A\cap B\cong S_3,\quad h^2=1,\quad
 \langle A,B\rangle=G_q,\quad AhA/A=\{sA:s\in S\}.
\]
The last set has four distinct cosets, and $S$ is stable under conjugation
by $A$.
\end{lemma}
\begin{proof}
On the trace-zero subspace with basis $(i,j,k)$, conjugation by the three
generators of $A_0$ and by $h_0$ gives, respectively,
\[
 \begin{pmatrix}1&0&0\\0&-1&0\\0&0&-1\end{pmatrix},\quad
 \begin{pmatrix}1&0&0\\0&0&-1\\0&1&0\end{pmatrix},\quad
 \begin{pmatrix}0&0&1\\1&0&0\\0&1&0\end{pmatrix},\quad
 H=\frac13\begin{pmatrix}-1&2&2\\2&-1&2\\2&2&-1\end{pmatrix}.
\]
The first three generate the $24$ determinant-one signed permutation
matrices. Their action on the four body diagonals identifies them with
$S_4$; they remain distinct modulo every odd prime. The adjoint action
is faithful projectively, also in odd characteristic. Conjugation by
$\alpha\in\mathbb H(\F_q)^\times$ is $L_\alpha R_{\alpha^{-1}}$ on
$\mathbb H(\F_q)$, of determinant $N(\alpha)^2N(\alpha)^{-2}=1$ for the
norm $N$; it fixes $1$ and preserves the norm form, so the adjoint action
embeds $G_q$ into $\mathrm{SO}_3(\F_q)$. The norm of $1+i$
is $2$, a nonsquare modulo $q$, so the determinant-square character is
non-trivial on $A$ and hence is its sign character. In particular the
transpositions of $A$ are exterior.

Here is a uniform intersection calculation. Let $K$ be the stabilizer in
$A$ of the line through $(1,1,1)$: its matrices are the positive even
permutation matrices and the negative odd permutation matrices. Thus
$K\cong S_3$, and every element of $K$ commutes with $H$.
Let $V$ be the four diagonal sign matrices of determinant one. Then
$A=VK$, and $V\cap K=1$. The three non-identity elements of $V$ are
conjugate under $K$. For $d=\operatorname{diag}(1,-1,-1)$,
\[
 HdH=\frac19\begin{pmatrix}-7&-4&-4\\-4&-1&8\\-4&8&-1\end{pmatrix}.
\]
Its second row has three nonzero entries modulo every prime $q\ge5$,
so $HdH\notin A$. Since $H$ commutes with $K$, the same conclusion
holds for the other two non-identity cosets of $K$ in $A$.
Therefore $A\cap HAH=K$, in characteristic zero and modulo every
allowed prime. This also proves $h\notin A$ and $B\ne A$.
The exterior $S_4$ is maximal in $G_q$ for these primes by
\cite[Theorem 3.5]{Giudici2007}; hence $\langle A,B\rangle=G_q$.

Finally $H^2=I$, and $A$ acts transitively by conjugation on the four
body diagonals $S$, with stabilizer $K$. The map
$A/K\longrightarrow AhA/A$, $aK\mapsto ahA$, is bijective by the
intersection calculation. Since $(aha^{-1})A=ahA$, its images are
exactly the four cosets represented by $S$.
\end{proof}

\begin{proposition}[The intersection-six stratum]
\label{prop:octahedral-cap6}
Let $A_1\ne A_2$ be exterior octahedral subgroups of $G_q$. Then
$|A_1\cap A_2|\le6$. If $|A_1\cap A_2|=6$, the ordered pair $(A_1,A_2)$
is simultaneously conjugate in $G_q$ to $(A,hAh^{-1})$, and for fixed
$A_1$ there are exactly four such $A_2$.
\end{proposition}
\begin{proof}
The exterior octahedral subgroups form one conjugacy class (Dickson's
classification of the subgroups of $\PSL_2(q)$ and
$\PGL_2(q)$~\cite{Huppert1967,King2005}; for the statement in this form
see~\cite[proof of Proposition~3.4]{Giudici2007}), so we may take $A_1=A$ and
$A_2=gAg^{-1}$ with $g\notin A$, in the orthogonal model of
\Cref{lem:octahedral-construction}.

\emph{Normalisers.} The three non-identity elements of $V$ have the
coordinate axes as fixed lines, so an element normalising $V$ permutes
these axes. Orthogonality forces its nonzero entries to be $\pm1$, and
determinant one leaves the $24$ elements of $A$. Thus $N_{G_q}(V)=A$;
since $V$ is characteristic in $A$, also $N_{G_q}(A)=A$.

\emph{Orders $8$ and $12$.} Put $D=A\cap A_2$. The determinant-square
character restricts to the sign character of $A$ and of $A_2$, so
$A\cap\PSL_2(q)$ and $A_2\cap\PSL_2(q)$ are their $A_4$ subgroups, whose
only subgroups of order $4$ are their normal Klein four subgroups. If
$|D|=12$, then $D$ is the $A_4$ of both groups. If $|D|=8$, then $D$
contains transpositions, so $D\cap\PSL_2(q)$ has order $4$ and is the
normal Klein four of both groups. In either case $A_2$ normalises $V$,
so $A_2\le A$, a contradiction. Hence $|D|\le6$, and $|D|=6$ forces
$D\cong S_3$.

\emph{Order $6$.} The four $S_3$ subgroups of $A$ are conjugate in $A$,
so replacing $g$ by $a_1ga_2^{-1}$ with $a_1,a_2\in A$, which does not
change $AgA$, we may assume that $g$ normalises $K$. Every automorphism of
$S_3$ is inner, so $N_{G_q}(K)=K\,C_{G_q}(K)$. The elements of $K$ are the
matrices $\operatorname{sgn}(\pi)P_\pi$, $\pi\in S_3$, so a matrix commutes
with $K$ exactly when it commutes with every permutation matrix $P_\pi$,
that is, when it is constant on the two orbits of $S_3$ on index pairs:
the commutant is $\{aI+bJ\}$, with $J$ the all-ones matrix. For
$M=aI+bJ$, orthogonality gives $a^2=1$ and $b(2a+3b)=0$; since $q\ne3$,
$M$ acts by $a+3b$ on $(1,1,1)$ and by $a$ on the complementary plane
$x+y+z=0$, so determinant one gives $a+3b=1$. The solutions are $I$ and
$H$. Hence $C_{G_q}(K)=\{1,h\}$, $N_{G_q}(K)=K\sqcup Kh$, and $g\in Kh$
because $g\notin A$. Therefore $AgA=AhA$; writing $g=ahb$ with $a,b\in A$
gives $A_2=a(hAh^{-1})a^{-1}$, so $(A,A_2)=a(A,hAh^{-1})a^{-1}$.

\emph{Count.} As $N_{G_q}(A)=A$, the map $gA\mapsto gAg^{-1}$ is
injective. The subgroups $A_2$ with $|A\cap A_2|=6$ are the images of the
cosets in $AhA/A$, of which there are $[A:K]=4$ by
\Cref{lem:octahedral-construction}.
\end{proof}

\subsection{Compression and the arithmetic spectral input}
Put $\mu=\frac14\sum_{s\in S}\delta_s$ and let $e_A$ denote the
uniform average over $A$ in $\C[G_q]$.

\begin{lemma}[The compression identity]
\label{lem:octahedral-compression}
In the group algebra,
\begin{equation}\label{eq:octahedral-compression}
 \mu e_A=e_A\mu=e_Ahe_A.
\end{equation}
Consequently $\rho(\mu)$ commutes with $P_A$ and its restriction to
$V_\rho^A$ is $P_A\rho(h)P_A$, whose norm equals $\opn{P_AP_B}$.
\end{lemma}
\begin{proof}
Conjugation invariance of $S$ gives commutation with $A$. Averaging
$ahe_A$ over $a\in A$ visits each of the four cosets in
\Cref{lem:octahedral-construction} equally often, proving
$e_Ahe_A=\frac14\sum_{s\in S}se_A$. Apply $\rho$ and
\Cref{cor:double-coset-angle}.
\end{proof}

\begin{proposition}[Quaternionic Ramanujan input]
\label{prop:octahedral-arithmetic}
For $\chi_q(g)$ the determinant-square character of $G_q$ and
$\epsilon_q=(3/q)$, the operator $\rho(\mu)$ has norm at most
$\sqrt3/2$ on every irreducible $\rho\notin\{\mathbf1,\chi_q\}$.
On those two characters its values are $1$ and $\epsilon_q$.
\end{proposition}
\begin{proof}
We apply the theorem of Davidoff--Sarnak--Valette
\cite[Theorem 4.2.2]{DSV2003}. Its hypotheses are distinct odd primes
$p,q$ with $q>2\sqrt p$, and its distinguished norm-$p$ generators
are specified in Section~2.6 of that reference. For $p=3$ they are
exactly the four classes of $\pm i\pm j\pm k$, so every prime here
satisfies the hypotheses. Remark~4.4.7 of the same reference derives the
Ramanujan bound from Eichler's theorem on weight-two cusp forms.

The vertex set of the walk follows from the construction. The set $S$ is
stable under conjugation by $A$ and contains $h$, so $\langle S\rangle$ is
normalised by $\langle A,h\rangle=G_q$ and is a non-trivial normal subgroup
of $\PGL_2(q)$. Since $\PSL_2(q)$ is simple with trivial centraliser in
$\PGL_2(q)$, $\langle S\rangle$ contains $\PSL_2(q)$. Every $s\in S$ has
determinant class $3$; hence $\langle S\rangle=\PSL_2(q)$ when
$\epsilon_q=1$ and $\langle S\rangle=\PGL_2(q)$ when $\epsilon_q=-1$. These
are the two vertex sets in the cited theorem, and the external input is
the spectral bound alone.

If $\epsilon_q=-1$, their Cayley graph on $G_q$ is connected,
bipartite and $4$-regular, with all eigenvalues other than $\pm4$
bounded in absolute value by $2\sqrt3$. The exceptional normalized
eigenspaces are the constants and $\chi_q$.
If $\epsilon_q=1$, the theorem gives a connected nonbipartite
Ramanujan graph on $\PSL_2(q)$. On $G_q$ the same walk is the union
of two copies, so the exceptional eigenspace at $1$ is spanned by
the constants and $\chi_q$, and there is no eigenvalue $-1$.
The regular representation contains every irreducible, giving the
assertion. In both cases each $s\in S$ has determinant square class $3$.
\end{proof}

The compression identity \eqref{eq:octahedral-compression} identifies the
Hecke operator. Left multiplication
by the fixed set $S$ on the cosets $G_q/A$ is generally a different
operator from the Hecke adjacency. For $q=5$, those two unnormalized
spectra are respectively $\{4,3,3,3,-1\}$ and
$\{4,-1,-1,-1,-1\}$. The latter graph is $K_5$.

\begin{theorem}[Sharp angle for the natural pair]
\label{thm:octahedral-sharp}
For the pair in \Cref{lem:octahedral-construction}, set
\[
 c_q=\max_{\rho\in\widehat{G_q},\,\rho\ne\mathbf1}
          \opn{P_A^\rho\rho(h)P_A^\rho}.
\]
Then, for every pair $A_1,A_2$ of exterior octahedral subgroups of $G_q$
with $|A_1\cap A_2|=6$,
\[
 \max_{\rho\ne\mathbf1}\opn{P_{A_1}^\rho P_{A_2}^\rho}
 =c_q=\max_{\rho\ne\mathbf1}\opn{P_A^\rho P_B^\rho}\le\frac{\sqrt3}{2},
 \qquad
 \lim_{q\to\infty}c_q
       =\frac{\sqrt3}{2}.
\]
Thus the constant is best possible uniformly over this family.
\end{theorem}
\begin{proof}
By \Cref{prop:octahedral-cap6}, $(A_1,A_2)=x(A,B)x^{-1}$ for some
$x\in G_q$, and $P_{xAx^{-1}}^\rho=\rho(x)P_A^\rho\rho(x)^{-1}$; this
gives the first equality in every representation. It remains to treat
$(A,B)$. For $\rho\notin\{\mathbf1,\chi_q\}$, the compression identity and
\Cref{prop:octahedral-arithmetic} give the upper bound. On $\chi_q$,
$P_A=0$ because $A$ is exterior, so this character contributes zero.

For the lower bound, define $X_q$ with vertices the left cosets $xA$,
and neighbours $xsA$ for $s\in S$. This is independent of the
representative $x$, since $S$ is invariant under conjugation by $A$.
The four neighbours are distinct and none is $xA$, by
\Cref{lem:octahedral-construction}. The relation is symmetric because
$h^{-1}=h$ and $AhA=(AhA)^{-1}$. The graph is connected because
$\langle A,h\rangle=G_q$, and has $\abs{G_q}/24$ vertices.
On functions on $G_q$ invariant under right multiplication by $A$,
its normalized adjacency is the restriction of the \emph{right}
convolution
\[
 (Wf)(x)=\frac14\sum_{s\in S}f(xs).
\]
Under the finite Peter--Weyl decomposition, the multiplicity spaces of
this coset representation are $V_\rho^A$ (or their duals, with identical
norms), and the blocks are the compressions
$P_A\rho(h)P_A$. Hence the absolute nonconstant spectral radius of
$X_q$ is $4c_q$. The trivial representation contributes just the
constant functions, and $\chi_q$ has no $A$-fixed vectors.
Alon--Boppana~\cite[Theorem 1.3.1]{DSV2003}, applied to these connected
loopless $4$-regular graphs with vertex count tending to infinity,
gives $\liminf 4c_q\ge2\sqrt3$. Together with the upper bound this
proves the limit.
\end{proof}

\subsection{From angles to tree blocks}
The minimum fixed-space angle $\theta_q=\arccos c_q$ satisfies
$\theta_q\ge\pi/6$ and $\theta_q\to\pi/6$ by
\Cref{thm:octahedral-sharp}. This angular separation controls how the
averages over $A$ and $B$ interact. The following estimate transfers it
to configurations whose sub-blocks approximate those averages.

\begin{lemma}[Quantitative gluing]
\label{lem:octahedral-gluing}
Let $P,Q$ be orthogonal projections with $\opn{PQ}\le c<1$. Suppose
$X=P+R$, $Y=Q+S$, where
$PR=RP=QS=SQ=0$ and $\opn R,\opn S\le\lambda<1$. Then
\begin{equation}\label{eq:octahedral-gluing}
 \opn{XYX^\dagger Y^\dagger}^2
 \le1-\frac{(1-\lambda)(1-c)}{8(1+\lambda)}.
\end{equation}
\end{lemma}
\begin{proof}
Put $g=1-\lambda^2$. Each factor $Z$ with associated projection $F$
satisfies
\[
 \norm w^2-\norm{Zw}^2\ge g\norm{(I-F)w}^2,
 \qquad \norm{w-Zw}\le(1+\lambda)\norm{(I-F)w}.
\]
Apply the four factors in their actual order to a unit $v$, and denote
their norm-squared losses by $e_j\ge0$. Their sum is
$E=1-\norm{XYX^\dagger Y^\dagger v}^2$.
At each step, telescoping back to $v$ and Cauchy--Schwarz give
\[
 \norm{(I-F)v}\le\frac{1+\lambda}{\sqrt g}
             \sum_{i\le j}\sqrt{e_i}
       \le\frac{2(1+\lambda)}{\sqrt g}\sqrt E.
\]
Both $P$ and $Q$ occur. Also $\opn{P+Q}\le1+c$: the map
$v\mapsto(Pv,Qv)$ has squared norm bounded by that of the block matrix
with identity diagonal and off-diagonal blocks $PQ,QP$.
Consequently
\[
 1-c\le\norm{(I-P)v}^2+\norm{(I-Q)v}^2
       \le\frac{8(1+\lambda)^2}{g}E.
\]
Taking the supremum over $v$ proves \eqref{eq:octahedral-gluing}.
\end{proof}

\begin{corollary}[Uniform contraction for natural-pair blocks]
\label{cor:octahedral-block}
Suppose a valid level-$3$ tree $\tau$ has level-$2$ leaf subgroups
$A,B\cong S_4$, exterior with $|A\cap B|=6$, and all leaves are exterior
involutions. For every non-trivial $\rho$,
\[
 \opn{M_\tau(\rho)}\le\kappa,
 \qquad
 \kappa=\left(1-\frac{183}{1736}
                       \left(1-\frac{\sqrt3}{2}\right)\right)^{1/2}<1.
\]
For a valid level-$m$ tree, $m\ge3$, all of whose level-$3$ sub-blocks
have this property, replace $\lambda$ by $\kappa$ and $D$ by $3$ in
\Cref{prop:quant}.
\end{corollary}
\begin{proof}
On the fixed space of its leaf subgroup, each level-$2$ operator is the
identity. On the orthogonal complement its norm is at most
$\lambda=17/200$: this is a finite rational certificate on the four
non-trivial irreducibles of $S_4$, over its $432$ valid ordered
transposition configurations generating $S_4$. The exact certificate
described in \Cref{sec:computation} constructs rational matrices with
invariant Gram matrices $G_\pi$ and verifies all principal minors of
$(17/200)^2G_\pi-D^tG_\pi D$ are nonnegative. Restriction of any ambient
unitary representation to $S_4$ decomposes into these irreducibles and
the trivial one; hence the certificate applies in every ambient $G_q$.
Both invariant subspaces are preserved, so the hypotheses of
\Cref{lem:octahedral-gluing} hold. Insert $c=\sqrt3/2$ from
\Cref{thm:octahedral-sharp}. The last assertion follows by expanding
into the level-$3$ blocks, as in \Cref{prop:quant}.
\end{proof}

The sharp angle estimate and the finite $S_4$ certificate thus give the
explicit contraction bound $\kappa$ for these octahedral blocks. Applying
\Cref{prop:stability} inside $S_4$ also gives a purely analytic, weaker
constant in the same corollary.

\subsection{Small primes and the remaining intersection orders}
\label{subsec:allpairs}
The values $c_q$ can be computed exactly for small $q$. Every
eigenspace of the graph $X_q$ in the proof of
\Cref{thm:octahedral-sharp} is invariant under left translation by
$G_q$, and each irreducible constituent of $\ell^2(G_q/A)$ has a
non-zero $A$-fixed vector; hence every eigenvalue of $X_q$ already
occurs on the functions on $A\backslash G_q/A$. Factoring the integer
characteristic polynomial of this quotient gives
\[
 c_{11}=\frac{1+\sqrt3}{4},\qquad c_{13}=\frac34,\qquad
 c_{19}=\frac{1+\sqrt5}{4}.
\]

\begin{table}[!ht]
\centering\small
\begin{tabular}{lrrrrr}
\toprule
$q$ & $11$ & $13$ & $19$ & $29$ & $37$ \\
\midrule
All-pairs maximum & $0.683013$ & $0.750000$ & $0.809017$ & $0.854361$ & $0.845651$ \\
\bottomrule
\end{tabular}
\smallskip
\caption{Numerical all-pairs maxima, to be compared with $\sqrt3/2\approx0.866025$.}
\label{tab:allpairs-evidence}
\end{table}

For pairs with $|A\cap B|<6$ no identification with an arithmetic
operator is available at present. \Cref{tab:allpairs-evidence} records,
for five primes, the maximum of $\opn{P_A^\rho P_{gAg^{-1}}^\rho}$ over
all $g\notin A$ and all non-trivial irreducible $\rho$, computed in
floating point from the operators induced on $A\backslash G_q/A$ by the
double cosets $AgA$. In each case the maximum is attained by the pairs
with $|A\cap gAg^{-1}|=6$ and agrees with $c_q$ to the precision shown.
Since an exterior octahedral subgroup is maximal, any two distinct ones
generate $G_q$, so these are the quantities in \Cref{conj:allpairs}.

\begin{remark}[Questions]
\label{rem:octahedral-questions}
By \Cref{prop:octahedral-cap6}, \Cref{conj:allpairs} concerns the pairs
with intersection of order $1$, $2$, $3$ or $4$. A weaker question is
whether the qualitative gap of \Cref{cor:almost-all-primes}, which holds
for all generating exterior pairs outside a set of primes of relative
density zero, holds for every prime $q\equiv\pm3\pmod8$.
\end{remark}


\section{Relation to word maps, conjugacy walks, and expansion}
\label{sec:related}
%% ===================================================================

The following connections identify the standard inputs and the questions addressed here.

\paragraph{Word maps.} Word maps connect random inputs in finite groups to
questions about images and distributions~\cite{LarsenShalev}.
For any fixed iterated commutator in distinct variables and any fixed
arrangement of brackets, Garion--Shalev~\cite{GarionShalev2009} prove
asymptotic equidistribution on finite non-abelian simple groups as their
orders tend to infinity. Their variables are uniform group elements, and
repeated occurrences of a variable use the same value.
Here each expanded occurrence contributes one independent fair bit:
the operator uses $XY\tilde X^{-1}\tilde Y^{-1}$ with independent copies
(\Cref{lem:tree,rem:notxy}). At fixed depth its output support has size
at most $2^{4^m}$, forcing depth $\Omega(\log\log|G|)$ for fixed-error
mixing. \Cref{ithm:E} attains this order in the specified families.

\paragraph{Conjugacy-class random walks.} When the leaf set is a single conjugacy class $C$,
uniform averaging over $C$ has eigenvalues $\chi(C)/\chi(1)$, with a
classical mixing theory~\cite{DS81,LulovPak}. In the commutator tree,
the leaf-subgroup criterion detects trapping through fixed spaces in
$\Ind_H^G\mathbf 1$. The non-identity target then forces escape from
proper solvable leaf subgroups at the depth given by \Cref{thm:dichotomy}.

\paragraph{Expansion and solvability.} Bourgain--Gamburd~\cite{BourgainGamburd}
and Salehi Golsefidy--Varj\'u~\cite{SalehiVarju} establish expansion under hypotheses
on generators, girth or the connected component of the Zariski closure, depending
on the theorem. In the tree model, a common fixed vector for a proper
leaf subgroup is the exact obstruction to strict contraction.
For bounded-depth involutive blocks, \Cref{cor:bounded-depth-gap} relates a
uniform gap to expansion of their particular leaf walks. Its upper comparison
uses Gao's projection union bound as presented in~\cite{AnshuAradVidick2016}.
The family rates of \Cref{ithm:E} use the bounded-rank diameter
theorem~\cite{PyberSzabo2016} and allow the fixed-depth gap to decrease
with group size. Strong uniform expansion gives the prime-field
refinement in \Cref{cor:almost-all-primes}.

\paragraph{Mixability of free subproducts.} Closest is the \emph{mixability} programme of Amir,
Blachar, Ghosh and Vishne \cite{AmirMixability}: when is a free random subproduct
$g_1^{\varepsilon_1}\cdots g_k^{\varepsilon_k}$ (arbitrary $g_i$, independent Bernoulli
$\varepsilon_i$) \emph{exactly} uniform? They characterise it---conjecturally---by $2'$-simplicity
(no non-trivial odd-order quotient), through a representation criterion
$\prod_i\bigl((1-p_i)I+p_i\rho(g_i)\bigr)=0$ whose singular-step lemma is a linear cousin of our
\Cref{lem:geom}. Their steps and Bernoulli parameters may be chosen freely. Here the expanded
steps are constrained by a balanced tree whose folded target must be non-trivial.
This validity condition brings the derived series into the threshold.
In particular, their mixability of every $S_n$ concerns different choices of
steps from our uniform statement over valid transposition trees.

\paragraph{Structural conclusion.} ``Mixing $\iff$ the steps generate $G$'' is classical for random walks.
For the independent-read commutator tree, the resulting structural statement is: the \emph{exact
uniform obstruction} is a proper non-solvable $\mathcal{C}$-generated subgroup, and the \emph{exact
level} is the largest derived length among proper $\mathcal{C}$-generated subgroups
(\Cref{thm:dichotomy}).

%% ===================================================================
\section{Computational verification}
\label{sec:computation}
%% ===================================================================

The artefact contains exact group computations, rational norm certificates,
and numerical checks. At level $2$, the tables enumerate reduced representatives:
interchanging leaves within either gate and interchanging the gates preserves
validity and the generated subgroup. Each valid orbit has eight ordered tuples.
These reductions certify validity and subgroup generation. Certified
numerical upper bounds on norms use exact checks covering all orientations;
NumPy singular values serve as cross-checks. The $S_5$ certificate is supplied
by the companion paper~\cite{ZermenoS5}.
The rank-one classification and the level-$5$ example are proved
algebraically. The script \path{psl33_maximals_check.py} independently
checks the orders, involution generation and derived lengths used in
\Cref{prop:psl33-level}.

\paragraph{The natural octahedral pair.}
\path{octahedral_pair_check.py} checks exact subgroup orders, intersections,
generation, coset adjacency and the compression identity for $q=5,11,13,19$;
it distinguishes the two coset operators exactly at $q=5$. Its larger
spectral checks are numerical. The supplementary Hamilton reduction check
uses exact integers. The script \path{s4_internal_gap_certify.py}
supplies the rational certificate used in \Cref{cor:octahedral-block},
checking all principal minors for the $432$ valid ordered
$S_4$-generating transposition configurations and all four non-trivial
irreducibles.
\path{octahedral_cap6_check.py} verifies, for $q=5,11,13,19,29,37$, the
intersection census of \Cref{prop:octahedral-cap6}, the uniqueness of the
double coset of intersection order six with its four neighbours, the
normaliser and centraliser computations, and the orthogonal dictionary.
It certifies the natural values $c_q$ exactly: the double-coset quotient
of the Hecke graph carries every distinct eigenvalue, and its integer
characteristic polynomial is factored over $\Z$ with rational root
isolation (SymPy). For $q=11,13,19$ this gives $c_q=(1+\sqrt3)/4$,
$3/4$, $(1+\sqrt5)/4$; the certificates for $q=29,37$ are recorded in the
artefact.
Arbitrary-order stability and the support bound
also have small-group numerical regression tests.

\paragraph{All-pairs numerical evidence.}
For an exterior $S_4$ subgroup $A<\PGL_2(q)$ and each non-trivial double
coset $AgA$, the non-trivial singular norm of the coset operator on
$\ell^2(G_q/A)$ equals that of its restriction to $A$-invariant
functions, an orbit-weighted matrix indexed by $A\backslash G_q/A$.
\path{octahedral_cap6_check.py} computes these norms by one
floating-point procedure for every prime in \Cref{tab:allpairs-evidence},
deflating a single copy of the trivial singular value, and cross-checks
them against dense singular value decompositions on $\ell^2(G_q/A)$ for
$q\le19$. The values are numerical. In each case the maximum is attained
at $|A\cap gAg^{-1}|=6$ and agrees, to the precision shown, with the
separately certified natural value $c_q$. The historical rank-one
experiments~\cite{ZermenoRankOneZenodo}, bundled with provenance as
\path{external/p2a/s4_allpairs_doublecoset_audit.py}, estimated the same
quantities by power iteration.

\begin{table}[htbp]
\centering\small
\begin{tabular}{lrrrrr}
\toprule
$q$ & $11$ & $13$ & $19$ & $29$ & $37$ \\
\midrule
All-pairs maximum & $0.683013$ & $0.750000$ & $0.809017$ & $0.854361$ & $0.845651$ \\
\bottomrule
\end{tabular}
\caption{Numerical all-pairs maxima, compared with $\sqrt3/2\approx0.866025$.}
\label{tab:allpairs-evidence}
\end{table}

\paragraph{Transvection levels.}
\path{transvection_levels.py} independently cross-checks the
$\SL_3(3)$ case of \Cref{prop:transvection-levels} by enumerating
transvection-generated subgroups up to conjugacy in exact arithmetic.
It finds $104$ transvections and eight non-trivial subgroup classes including $G$;
the largest proper derived length is $4$, attained at order $216$.

\paragraph{The generation criterion (reduced level-$2$ universe).}
For each reduced valid configuration we compare the numerical regular-representation
second singular value with its exact leaf-subgroup closure, as predicted by
\Cref{thm:criterion}. The counts in \Cref{tab:criterion} are reduced; the respective
ordered totals are $105168$ and $219360$.

\begin{table}[t]
\centering\small
\begin{tabular}{lrrrrr}
\toprule
Group & Reduced & Intransitive & Proper & Generate $G$ & Mismatches \\
\midrule
$\GL_3(\F_2)$ & $13146$ & $2772$ & $2772$ & $10374$ & $0$ \\
$S_5$ on $6$ pts & $27420$ & $0$ & $4695$ & $22725$ & $0$ \\
\bottomrule
\end{tabular}
\caption{Reduced level-$2$ all-involution checks. ``Intransitive'' counts augmentation norm one;
``Proper'' counts regular second singular value one, agreeing with the exact leaf
closure in every case. The augmentation test misses $4695$ proper transitive
closures for $S_5$ on six points ($3705$ of order $60$, $990$ of order $24$).}
\label{tab:criterion}
\end{table}

\paragraph{Mixing levels.}
\Cref{tab:levels} lists the level at which contraction sets in, with its certification
status. The $S_5$ five-point row reproduces the companion paper~\cite{ZermenoS5} to the digit
for the reduced enumeration ($\lambda\approx0.0353047532$ over $73810$ reduced
level-$3$ configurations). The corrected full oriented enumeration has maximum
$\lambda\approx0.03540023208379892$, so the bound $71/2000=0.0355$ is certified in exact integer arithmetic by
\path{20_level3_exact_certificate.py} from~\cite{ZermenoS5}, supplied unchanged in the artefact
with its provenance. It covers $9447840$ oriented pairs through $39591$ symmetry
orbits; the sign representation vanishes separately.

\begin{table}[t]
\centering\small
\begin{tabular}{p{.33\textwidth}p{.22\textwidth}p{.34\textwidth}}
\toprule
Action and leaves & Stabiliser / dl & Level and quantitative evidence \\
\midrule
$S_5$, $5$ points, transpositions & $S_4$ / $3$ & $3$; exact bound $0.0355$~\cite{ZermenoS5} \\
$S_5$, $6$ points, all involutions & $F_{20}$ / $2$ & $2$ in augmentation only; full contraction obstructed \\
$\PGL_2(7)$, $8$ points, exterior & Borel / $2$ & $2$ analytically; norms sampled \\
$\GL_3(\F_2)$, $7$ points, involutions & $S_4$ / $3$ & $3$ analytically; level-$3$ norms capped \\
$S_6$, $6$ points, transpositions & $S_5$ / non-solvable & No uniform level; obstructed configurations persist \\
\bottomrule
\end{tabular}
\caption{Contraction levels and quantitative evidence.}
\label{tab:levels}
\end{table}

\paragraph{A degenerate example.}
$\GL_2(\Z_4)$ produces zero valid level-$3$ configurations,
confirming the degeneracy of \Cref{sec:layers}.

%% ===================================================================
\section{Open problems}
\label{sec:open}
%% ===================================================================

\begin{enumerate}[nosep,leftmargin=1.4em]
\item \textbf{Uniform expansion and sharp angles.}
  For prime $q\equiv\pm3\pmod8$, $q\ne3$, is
  \[
   \sup_{\substack{A,B\cong S_4\ \mathrm{exterior}\\
                    \langle A,B\rangle=\PGL_2(q)}}\quad
   \sup_{\rho\ne\mathbf1}\opn{P_A^\rho P_B^\rho}
   \le\sqrt3/2+o(1)\qquad(q\to\infty)?
  \]
  This question was formulated in the preliminary rank-one
  preprint~\cite[Conjecture 6.9]{ZermenoRankOneZenodo}.
  \Cref{prop:octahedral-cap6,thm:octahedral-sharp} prove it for every
  pair with $|A\cap B|=6$, so it remains open for the intersection
  orders $1,2,3,4$; \Cref{tab:allpairs-evidence} records finite
  numerical evidence.
  A separate question is whether the qualitative gaps of
  \Cref{cor:almost-all-primes} for all exterior generating pairs and
  all ambient-generating level-$2$ involutive blocks extend to every
  admissible prime.
  Finally, do the involution classes of the minimal simple families
  $\PSL_2(2^\ell)$, $\PSL_2(3^\ell)$ and $\mathrm{Sz}(2^\ell)$
  admit a constant gap for every valid level-$D$ block, uniformly as
  the prime exponent $\ell$ grows (odd in the latter two families)?
  \Cref{cor:few-exceptional-blocks} allows exceptional configurations;
  removing them in these families is a concrete consequence sought
  from the bounded-rank uniform expansion conjecture discussed in
  \cite[Introduction]{BeckerBreuillard2026}.
\item \textbf{Fixed templates.} Characterise which \emph{prescribed} rigid families attain the
  configurable bound $D$: \Cref{prop:rigidupper} gives $L(F)\le D$ for all of them,
  \Cref{thm:s5levels} shows that over $S_5$ both $D$ and $D-1$ occur, and
  \Cref{cor:binarylevel} characterises prefixes admitting an exact-$D$
  continuation. For a prescribed family, determine the optimal contraction
  rate and whether its base-level gap stays positive as the ambient group
  varies. Submultiplicativity propagates a given base gap; the sharp rate
  remains to be determined.
\item \textbf{Irreducible matrix classes and congruence levels.}
  For $n\ge3$, classify the generating conjugacy classes of
  $\GL_n(\F_q)$ with irreducible characteristic polynomial for which
  every proper class-generated subgroup is solvable. These elements
  preserve no proper nonzero subspace, so this question concerns the
  non-parabolic maximal subgroups. \Cref{prop:transvection-levels}
  settles the contrasting transvection family in $\SL_n(q)$.
  For a generating conjugacy class of $\GL_2(\Z/p^2)$, $p\ge5$,
  whose reduction is unobstructed, determine exactly when the level
  in \Cref{prop:congruence-class} is $\bar D$ or $\bar D+1$.
  More generally, characterise the proper class-generated supplements
  to congruence kernels when unions of classes are allowed, starting
  from \Cref{rem:congruence-supplement}.
\end{enumerate}

\section*{AI Declaration}
% Before submission, complete the model/version inventory for each tool from
% session records or the author's records; tool names alone do not identify models.
OpenAI ChatGPT/Codex, Anthropic Claude and Google Gemini were used to
search for counterexamples and errors in proofs, assist with developing
and verifying code, and edit the language.
The author is responsible for the mathematical statements, proofs,
references, code and exposition.

\section*{Acknowledgements and declarations}
\paragraph{Mathematical software.}
Python and NumPy were used for the computations described in
\Cref{sec:computation}; SymPy was used to factor integer characteristic
polynomials and isolate roots in the exact octahedral certificates.

\paragraph{Funding.}
This research received no specific grant from any funding agency in the public, commercial or
not-for-profit sectors.

\paragraph{Conflict of interest.}
The author declares no conflict of interest.

\paragraph{Data access.}
The verification artefact and its sources are publicly deposited at
\url{https://github.com/ernest-zermeno/commutator-tree-mixing-finite-groups}.
Its README and manifest identify exact certificates, numerical checks, sampled
runs and the certificate from~\cite{ZermenoS5}.

For \Cref{subsec:rigid}, \path{bridge_realizability_certificate.py}
checks the $S_5$ rigid families, \path{target_transition_graph.py}
checks target transitions, and \path{binary_target_grammar_certificate.py}
checks the binary persistence criterion in the groups listed in its documentation.

\bibliographystyle{plainnat}
\bibliography{refs}

\end{document}
